% 第 8 章 PXL（前半：p175--210）
\bichapter{PXL}{PXL}
\label{chap:pxl}

本章介绍可编程可扩展语言 PXL，以及在 IC Validator 工具中函数定义与函数使用的基础。

PXL 包括以下各节：
\begin{itemize}
  \item PXL 体系结构（PXL Architecture）
  \item 一般定义（General Definitions）
  \item 变量与运算符（Variables and Operators）
  \item 数据类型与类型系统（Data Types and Typing）
  \item 流程控制（Flow Control）
  \item 函数（Functions）
  \item 程序结构（Program Structure）
  \item 模块化代码示例（Example of Modularized Code）
\end{itemize}

% ============================================================
\section{PXL 体系结构}
\subsection*{PXL Architecture}

PXL 完全可编程，具备下列一般可编程性特性：
\begin{itemize}
  \item 变量（Variables）
  \item 流程控制（Flow Control）
  \item 带标准索引的列表（Lists / arrays）
  \item 结构体（Structures）
  \item 哈希表（Hashes）
  \item 枚举类型（Enumerated types）
  \item 用户定义函数（User-defined functions）
  \item 数学、逻辑与字符串运算
  \item 支持命令行选项的宏（Macros）
  \item 环境变量（Environment variables）
  \item 静态作用域（Static scope）
\end{itemize}

这些特性可带来：
\begin{itemize}
  \item 直观的 runset 开发
  \item 易于维护的 runset
  \item 简洁的代码
  \item 通过子程序实现代码复用
  \item 支持用户定义函数
\end{itemize}

图~\ref{fig:pxl-structure} 给出 PXL 的基本结构。

\figplaceholder{Figure 55: PXL Structure}{PXL 结构}{fig:pxl-structure}

% ============================================================
\section{一般定义}
\subsection*{General Definitions}

本节给出 PXL 中部分语法元素的定义与特征，例如标识符、字符串字面量、关键字、运算符与注释。理解这些语法元素的用途与作用，有助于编写 runset。

\subsection{标识符}
\subsubsection*{Identifiers}

任意 PXL 组件的标识符均可按任意顺序包含字母字符、数字与下划线。
PXL 区分大小写。下列标识符彼此不同：
\begin{itemize}
  \item \cmd{abCD}
  \item \cmd{ABCD}
  \item \cmd{abcd}
\end{itemize}

标识符须遵守下列命名限制：
\begin{itemize}
  \item 双下划线前缀（\cmd{\_\_}）。不允许以两个下划线开头的标识符；保留给语言内部使用。
  \item 单下划线前缀（\cmd{\_}）。不允许以一个下划线开头的标识符；保留给产品使用。
  \item 全大写。保留给应用程序编程接口（API）使用。不过，可将全大写标识符用作枚举类型元素。
  \item 保留关键字。不得将保留关键字用作标识符。PXL 保留关键字列表见第~\ref{sec:pxl-keywords}~节「关键字」。
\end{itemize}

下列名称也不可用作标识符：
\begin{itemize}
  \item 函数名：这些函数定义于 \textit{IC Validator Reference Manual} 与 IC Validator 头文件中。示例函数名包括 \cmd{capacitor}、\cmd{check\_property}、\cmd{internal1} 与 \cmd{xor}。
  \item 容器类型名：列表、哈希与结构体的名称定义于 IC Validator 头文件中。容器类型名示例：
  \begin{itemize}
    \item 列表：\cmd{coordinate\_l} 用于定义 \cmd{property()} 函数
    \item 哈希：\cmd{layer\_list\_h} 用于定义 \cmd{density()} 函数
    \item 结构体：\cmd{check\_property\_s} 用于定义 \cmd{check\_property()} 函数
  \end{itemize}
\end{itemize}

合法标识符示例：
\begin{lstlisting}[style=pxl]
This
IsAnIdentifier
another_identifier_using_underscores
layer20
\end{lstlisting}

保留（非法）标识符示例：
\begin{lstlisting}[style=pxl]
__SYSTEMNAME
\end{lstlisting}
（双下划线前缀；保留给语言内部使用）

\begin{lstlisting}[style=pxl]
_productname
\end{lstlisting}
（单下划线前缀；保留给产品使用）

\begin{lstlisting}[style=pxl]
HIERARCHICAL
\end{lstlisting}
（保留关键字）

\subsection{字符串字面量}
\subsubsection*{String Literals}

字符串字面量是以双引号（\cmd{"}）开始并结束的一组字符。

\begin{noteBox}
文本字符串的第一个字符不能是整数，也不能是后跟整数的空格。\\
字符串字面量不得包含换行。若字符串字面量中需要换行，必须使用换行字符（\cmd{\textbackslash n}）。
\end{noteBox}

下列字符组均为字符串字面量示例：
\begin{lstlisting}[style=pxl]
"hello"
"\""
"a string with\nmultiple lines"
"abc" "def"
\end{lstlisting}

字符串字面量会自动与相邻的字符串字面量合并（字符串拼接）。尤其是，它们会与包含字符串值的预处理器变量（如 \cmd{\_\_FILE\_\_}）以及环境变量引用（如 \cmd{\$<id>}）合并。

可在字符串内使用转义序列，将部分字符串内容标记为给编译器的指令。PXL 可用的转义序列见表~\ref{tab:pxl-escape}。

\begin{longtable}{@{}>{\raggedright\arraybackslash\ttfamily}p{0.22\textwidth} p{0.70\textwidth}@{}}
\caption{转义序列（Escape Sequences）}\label{tab:pxl-escape}\\
\toprule
\textbf{\textrm{转义}} & \textbf{\textrm{定义}} \\
\midrule
\endfirsthead
\toprule
\textbf{\textrm{转义}} & \textbf{\textrm{定义}} \\
\midrule
\endhead
\bottomrule
\endfoot
\textbackslash n & 换行字符 \\
\textbackslash" & 双引号字符 \\
\textbackslash\textbackslash & 单个反斜杠字符 \\
\textbackslash NNN & 字符的八进制值 \\
\textbackslash t & 制表符 \\
\textbackslash xNN & 字符的十六进制值 \\
\end{longtable}

\subsection{数值常量}
\subsubsection*{Numeric Constants}

PXL 中的数值常量定义见表~\ref{tab:pxl-numeric}。

\begin{longtable}{@{}p{0.26\textwidth} p{0.42\textwidth} >{\raggedright\arraybackslash\ttfamily}p{0.24\textwidth}@{}}
\caption{数值常量（Numeric Constants）}\label{tab:pxl-numeric}\\
\toprule
\textbf{数值类型} & \textbf{条件} & \textbf{\textrm{示例}} \\
\midrule
\endfirsthead
\toprule
\textbf{数值类型} & \textbf{条件} & \textbf{\textrm{示例}} \\
\midrule
\endhead
\bottomrule
\endfoot
double & 含指数 & 4e5, 10.305e+10 \\
double & 含小数点 & 3.141592 \\
十进制整数（decimal integer） & 以数字 1--9 开头 & 1023 \\
八进制数（octal number） & 以零（0）开头，且各位为 0--7 & 0346 \\
十六进制数（hexadecimal number） & 以 \cmd{0x} 或 \cmd{0X} 开头，且仅含十六进制数字（0--9、a--f、A--F） & 0xDEADBEEF \\
\end{longtable}

\subsection{关键字}
\subsubsection*{Keywords}
\label{sec:pxl-keywords}

关键字是语言识别的内置词。关键字为保留字，因此不能用作标识符。

关键字的若干用途包括：
\begin{itemize}
  \item 流程控制，如 \cmd{if}、\cmd{foreach} 与 \cmd{while}
  \item 运算符，如 \cmd{in} 与 \cmd{contains}
  \item 常量，如 \cmd{true} 与 \cmd{false}
\end{itemize}

下列各表列出 PXL 中的保留字。

\begin{noteBox}
有关命名限制的更多信息，见「标识符」小节。
\end{noteBox}

表~\ref{tab:pxl-sys-kw} 列出由系统保留、外部不可用的关键字。

\begin{longtable}{@{}>{\raggedright\arraybackslash\ttfamily}p{0.30\textwidth} >{\raggedright\arraybackslash\ttfamily}p{0.30\textwidth} >{\raggedright\arraybackslash\ttfamily}p{0.30\textwidth}@{}}
\caption{系统关键字（System Keywords）}\label{tab:pxl-sys-kw}\\
\toprule
\endfirsthead
\endhead
\bottomrule
\endfoot
barrier & broadcast & builtin \\
deprecated & intrinsic & method \\
obsolete & preprocess & primary \\
proxy\_out & published & required \\
shutdown & unsafe & unique \\
volatile &  &  \\
\end{longtable}

表~\ref{tab:pxl-indesign-kw} 列出 In-Design 流程使用的关键字。请勿在 IC Validator runset 中使用这些名称作为变量。

\begin{longtable}{@{}>{\raggedright\arraybackslash\ttfamily}p{0.48\textwidth} >{\raggedright\arraybackslash\ttfamily}p{0.48\textwidth}@{}}
\caption{In-Design 流程关键字（In-Design Flow Keywords）}\label{tab:pxl-indesign-kw}\\
\toprule
\endfirsthead
\endhead
\bottomrule
\endfoot
INCREMENTAL\_MILKYWAY\_OPTIONS\_INCLUDE\_VIEW & INCREMENTAL\_OPTIONS\_SELECT\_WINDOW \\
INDESIGN\_DO\_FLOATING\_VIA\_FILL\_PURGE & INDESIGN\_DONT\_UPDATE\_TECHNOLOGY\_FILE \\
INDESIGN\_ENABLE\_MY\_ASSIGN & METAL\_FILL\_BLOCKAGE\_LAYERS \\
METAL\_FILL\_BLOCKAGE\_VIEW & METAL\_FILL\_DATATYPES \\
METAL\_FILL\_DENSITY\_DEBUG & METAL\_FILL\_DENSITY\_CHECK \\
METAL\_FILL\_DENSITY\_DELTA\_WIN\_SIZE & METAL\_FILL\_DENSITY\_MIN \\
METAL\_FILL\_FLATTEN & METAL\_FILL\_INLIB \\
METAL\_FILL\_INLIB\_CELL & METAL\_FILL\_INLIB\_EXCLUDE\_FILL \\
METAL\_FILL\_INLIB\_PATH & METAL\_FILL\_INSTALL\_ICV\_RH \\
METAL\_FILL\_INSTALL\_PRIMEYIELD\_RH & METAL\_FILL\_LAYER\_NUMBERS \\
METAL\_FILL\_OUTLIB & METAL\_FILL\_OUTLIB\_APPLY\_PREFIX \\
METAL\_FILL\_OUTLIB\_CELL & METAL\_FILL\_OUTLIB\_CELL\_PREFIX \\
METAL\_FILL\_OUTLIB\_MODE & METAL\_FILL\_OUTLIB\_PATH \\
METAL\_FILL\_OUTLIB\_TEMP\_CELL & METAL\_FILL\_OUTLIB\_VIEW \\
METAL\_FILL\_OUTLIB\_RF\_CELL\_PREFIX & METAL\_FILL\_RUNSET2LIBRARY\_LAYER\_MAP \\
METAL\_FILL\_SELECT\_WINDOW & METAL\_FILL\_USER\_RUNSET \\
\_METHODOLOGY\_FUNCTIONS\_RS\_ & SAVE\_METAL\_FILL\_INLIB\_EXCLUDE\_FILL \\
\_SMF\_ICV\_RH\_ & SNPSINDESIGN \\
VIA\_FILL\_ALLOWED\_ONE\_SIDE & VIA\_FILL\_ASSOCIATED\_METAL \\
VIA\_FILL\_DATATYPES & VIA\_FILL\_LAYER\_NUMBERS \\
VIA\_FILL\_MASK\_NAMES & VIA\_FILL\_MET\_ENCLOSURE \\
VIA\_FILL\_OUTLIB\_RF\_CELL\_PREFIX & \_WRITE\_AREFS \\
\end{longtable}

表~\ref{tab:pxl-app-ids} 列出不可用作标识符的保留应用程序标识符。

\begin{longtable}{@{}>{\raggedright\arraybackslash\ttfamily}p{0.30\textwidth} >{\raggedright\arraybackslash\ttfamily}p{0.30\textwidth} >{\raggedright\arraybackslash\ttfamily}p{0.30\textwidth}@{}}
\caption{保留应用程序标识符（Reserved Application Identifiers）}\label{tab:pxl-app-ids}\\
\toprule
\endfirsthead
\endhead
\bottomrule
\endfoot
binary & boolean & break \\
by & const & constraint \\
constraint\_category & CONSTRAINT\_EQ & CONSTRAINT\_GE \\
CONSTRAINT\_GELE & CONSTRAINT\_GELT & CONSTRAINT\_GT \\
CONSTRAINT\_GTLE & CONSTRAINT\_GTLT & CONSTRAINT\_LE \\
CONSTRAINT\_LT & CONSTRAINT\_NE & continue \\
DELETE\_TEXT & DELETE\_AND\_REPORT & defined \\
double & elif & else \\
elseif & enum & false \\
for & foreach & function \\
hash & if & in \\
in\_out & integer & list \\
newtype & of & on \\
out & return & returning \\
step & string & struct \\
substrate & SUBSTRATE & thru \\
to & true & unary \\
undefined & violation & void \\
while &  &  \\
\end{longtable}

枚举类型元素也不能用作标识符。表~\ref{tab:pxl-enum-elems} 列出 IC Validator 工具使用的枚举类型元素。

\begin{longtable}{@{}>{\raggedright\arraybackslash\ttfamily}p{0.30\textwidth} >{\raggedright\arraybackslash\ttfamily}p{0.30\textwidth} >{\raggedright\arraybackslash\ttfamily}p{0.30\textwidth}@{}}
\caption{枚举类型元素（Enumerated Type Elements）}\label{tab:pxl-enum-elems}\\
\toprule
\endfirsthead
\endhead
\bottomrule
\endfoot
ABORT & ABSOLUTE & ABSTRACT\_DESIGN \\
ACCEPT & ACUTE & ADJUST \\
ALIEN & ALL & ALL\_CELLS \\
ALL\_POLYGONS & ALLOW & ALWAYS \\
AMBIT\_ON\_BLOCK & ANALOG\_DESIGN & ANALOG\_GROUND\_NET \\
ANALOG\_POWER\_NET & ANALOG\_SIGNAL\_NET & ANNOTATION\_FILE \\
ANTENNA\_CELL & ANTIDISHING\_BRIDGE & ANY\_DEVICE\_NAME \\
ANY\_LEVEL & APPEND & APPEND\_NEW\_VERSION \\
AREA\_FILL & AREA\_FILL\_USE & ARC \\
ARC\_PLAN\_CELL & ASSIGNS & ASYMMETRIC \\
ASYMMETRIC\_INTERSECTING & AUTO & AUTO\_DETECT \\
AUTOMATIC & AXIS & BACKSLASH \\
BALANCED & BASE & BISECTOR \\
BLACK\_BOX\_DESIGN & BLOCK\_RADIAL & BLOCK\_SQUARE \\
BLOCKAGE\_OBJECT & BODY\_EXTENTS & BODY\_POLYGONS \\
BOTH & BOUNDARY & BOTTOM \\
BOTTOM\_HEIGHT & BREAK\_PATH & BRIEF \\
BULK & BUS & BV\_CELL \\
CAPACITOR & CASE\_INSENSITIVE\_NAME & CATEGORY\_BLOCKAGE \\
CDL & CELL\_LEVEL & CELL\_ROW \\
CENTER & CENTER\_TO\_CENTER & CENTERLINE \\
CHECK & CHIP & CHIP\_CONTEXT \\
CHIP\_CELL & CHIPWIDE\_AVERAGE & CLASSIFY \\
CLEAN & CLIP & CLIPPED\_DELTA\_WINDOW \\
CLIPPED\_CENTER & CLOCK & CLOCK\_0\_SKEW \\
CLOCK\_AND\_SPLIT\_CLOCK & CLOCK\_BUFFER\_CELL & CLOCK\_NET \\
CLOCK\_RING & CLOCK\_STRIPE & CLOSED\_OUTER\_BOUNDARY \\
CLOSEST & CMPSIM & COINCIDENT \\
COLLECTOR & COLOR\_1 & COLOR\_2 \\
COLOR\_LAYER\_TYPE & COLOR\_LINE\_END & COLOR\_RING \\
COLOR\_SIDE & COMBINE\_COLOR & COMMENT \\
CONCAVE\_TO\_CONCAVE & CONCAVE\_TO\_EDGE & CONDUCTING\_LAYER\_TYPE \\
CONDUCTIVE\_HEIGHT & CONNECT & CONFLICT \\
CONTACT & CONTACTARRAY & CONTINUE \\
CONVERT\_TO\_POLYGON & CONVERT\_TO\_RECT & CONVEX\_TO\_CONCAVE \\
CONVEX\_TO\_CONVEX & CONVEX\_TO\_EDGE & COORDINATES\_AND\_DEVICE\_TYPES \\
CORE\_AREA\_OBJECT & CORE\_OF\_BLOCK & CORE\_WIRE \\
CORE\_WIRE\_USE & CORNER & CORNER\_PAD\_CELL \\
CORNER\_PAD\_DESIGN & CORNER\_TO\_CORNER & CORNER\_TO\_CORNER\_OR\_EDGE \\
CORNER\_TO\_EDGE & COUNT & COVER\_CELL \\
COVER\_DESIGN & CREATED\_PORTS & CROSS \\
CUMULATIVE\_TOPOGRAPHY & CUSTOM & CUTTING \\
DATATYPE\_ZERO\_ONLY & DEEP\_NWELL\_NET & DEEP\_PWELL\_NET \\
DEFOCUS & DELTA\_WINDOW & DENSITY \\
DENSITY\_GRADIENT & DENSITY\_NORMAL & DENSITY\_PERIMETER \\
DENSITY\_SMARTSTEP & DENSITY\_TOTARGE & DENSITY\_UNDER\_1 \\
DENSITY\_UNDER\_2 & DEPTH\_OF\_FOCUS\_COLUMNS & DEPTH\_OF\_FOCUS\_COLUMNS\_HIGH \\
DEPTH\_OF\_FOCUS\_COLUMNS\_LOW & DEPTH\_OF\_FOCUS\_ROWS & DEPTH\_OF\_FOCUS\_ROWS\_HIGH \\
DEPTH\_OF\_FOCUS\_ROWS\_LOW & DESIGN & DESIGN\_VIEW \\
DETAIL\_ROUTE & DETAIL\_ROUTE\_USE & DEVICE\_NAME \\
DEVICE\_PIN & DIAMOND & DIELECTRIC\_HEIGHT \\
DIFFERENT\_NET & DIFFERENT\_POLYGON & DIFFERENT\_SIZE\_RECTANGLES \\
DIODE\_DESIGN & DISCARD & DONUT \\
DOT & DOUBLE & DOUBLE\_HEIGHT\_STANDARD\_CELL \\
DOUBLE\_LIST & DRC & DROP \\
DUAL & DYNAMIC\_SHIELD & EDGE \\
EDGES & EDGE\_TO\_CONCAVE & EDGE\_TO\_CONVEX \\
EDIF & EIGHT & ELLIPSE \\
ELLIPTICAL & ELPC & EMITTER \\
EMPTY & ENCLOSE & END\_CAP\_DESIGN \\
EQUATE\_NETS & ERROR & ERRORS \\
ERROR\_CLASSIFICATION & ERROR\_NO\_ABORT & EXCLUSIVE \\
EXPANDED\_EDGE & EXPLODE & EXTENDED \\
EXTENTS & EXTERIOR & EXTERIOR\_INTERIOR \\
EXTRACT & FAIL & FAILED \\
FAILED\_INSERTION & FC\_BUMP\_CELL & FC\_DRIVER\_CELL \\
FEEDTHRU\_BLOCKAGE & FEMTO & FILL \\
FILL\_BLOCKAGE & FILL\_BLOCKAGE\_USE & FILL\_DESIGN \\
FILL\_RINGS & FILL\_TRACK & FILL\_USE \\
FILLER\_CELL & FILLER\_DESIGN & FIRST\_PATTERN\_LAYER \\
FIXED & FLIP\_CHIP & FLIP\_CHIP\_DRIVER\_DESIGN \\
FLIP\_CHIP\_PAD\_CELL & FLIP\_CHIP\_PAD\_DESIGN & FLIP\_FLOP\_CELL \\
FLOATING & FOLLOW\_PIN & FOLLOW\_PIN\_USE \\
FORTY\_FIVE & FOUR & FRAM \\
FRAME\_VIEW & FULL & FULL\_NAME \\
FULL\_NAME\_CASE\_INSENSITIVE & FULL\_NAME\_CASE\_SENSITIVE & GALAXY\_AND\_XO\_CELL \\
GALAXY\_CELL & GDSII & GENERIC \\
GLOBAL & GLOBAL\_REPLACE & GLOBAL\_ROUTE \\
GLOBAL\_ROUTE\_USE & GROUND & GROUND\_NET \\
GROW & GUIDANCE & HALF\_WIDTH \\
HIERARCHICAL & HIGH & HIGHEST\_TEXT \\
HOLDING\_CELL & HORIZONTAL & HORIZONTAL\_AXIS \\
HORIZONTAL\_TRACK & HORIZONTAL\_TRACK\_OBJECT & HORIZONTAL\_WIRE\_TRACK \\
HORIZONTALWIRE & HYBRID & HYBRID\_BAIL\_EARLY \\
HYBRID\_PATTERN\_ONLY & ICV & ICV\_LCC\_AOI \\
ICV\_LCC\_BLOCKING & ICV\_LCC\_CO & ICV\_LCC\_CONTOUR\_AREA \\
ICV\_LCC\_M1 & ICV\_LCC\_M1\_DUMMY & ICV\_LCC\_M1\_OPC\_DUMMY \\
ICV\_LCC\_M2 & ICV\_LCC\_M2\_DUMMY & ICV\_LCC\_M2\_OPC\_DUMMY \\
ICV\_LCC\_M3 & ICV\_LCC\_M3\_DUMMY & ICV\_LCC\_M3\_OPC\_DUMMY \\
ICV\_LCC\_M4 & ICV\_LCC\_M4\_DUMMY & ICV\_LCC\_M4\_OPC\_DUMMY \\
ICV\_LCC\_M5 & ICV\_LCC\_M5\_DUMMY & ICV\_LCC\_M5\_OPC\_DUMMY \\
ICV\_LCC\_M6 & ICV\_LCC\_M6\_DUMMY & ICV\_LCC\_M6\_OPC\_DUMMY \\
ICV\_LCC\_M7 & ICV\_LCC\_M7\_DUMMY & ICV\_LCC\_M7\_OPC\_DUMMY \\
ICV\_LCC\_M8 & ICV\_LCC\_M8\_DUMMY & ICV\_LCC\_M8\_OPC\_DUMMY \\
ICV\_LCC\_M9 & ICV\_LCC\_M9\_DUMMY & ICV\_LCC\_M9\_OPC\_DUMMY \\
ICV\_LCC\_M10 & ICV\_LCC\_M10\_DUMMY & ICV\_LCC\_M10\_OPC\_DUMMY \\
ICV\_LCC\_MAX & ICV\_LCC\_OD & ICV\_LCC\_OD\_DUMMY \\
ICV\_LCC\_OD\_OPC\_DUMMY & ICV\_LCC\_PO & ICV\_LCC\_PO\_DUMMY \\
ICV\_LCC\_PO\_OPC\_DUMMY & ICV\_LCC\_VIA1 & ICV\_LCC\_VIA2 \\
ICV\_LCC\_VIA3 & ICV\_LCC\_VIA4 & ICV\_LCC\_VIA5 \\
ICV\_LCC\_VIA6 & ICV\_LCC\_VIA7 & ICV\_LCC\_VIA8 \\
ICV\_LCC\_VIA9 & IF\_INSERTION\_FAILS & IGNORE \\
IGNORE\_CAP\_LAYER\_TYPE & IMAGE\_CELL & IN \\
INCLUSIVE & INCLUSIVE\_PARALLEL & INDIVIDUAL \\
INDUCTOR & INNER & INPUT\_LAYOUT \\
INSERTION & INSIDE & INSIDE\_TO\_OUTSIDE \\
INSTANCE\_NAME & INSTANCES & INTERACT \\
INTERACTING & INTERIOR & INTERNAL\_DESIGNS \\
INTERSECTION & IO\_PAD\_CELL & KEEP \\
KEEP\_ALL & KEEP\_ENCLOSED & KEEP\_INTERACTING \\
KEEP\_ONE & KILO & LATCH\_CELL \\
LAYER & LAYER\_INTENT\_BASE & LAYER\_INTENT\_DEVICE \\
LAYER\_INTENT\_FILL & LAYER\_INTENT\_INTERCONNECT & LAYER\_INTENT\_METAL \\
LAYER\_INTENT\_VIA & LAYER\_MAPS & LAYER1 \\
LAYER2 & LAYOUT\_UNTEXTED & LAYOUT\_VIEW \\
LCC & LEVEL\_SOURCE\_DRAIN & LI\_BASE \\
LI\_FILL & LI\_INTERCONNECT & LI\_METAL \\
LI\_VIA & LIB\_CELL\_DESIGN & LIB\_CELL\_PIN\_CONNECT \\
LIB\_CELL\_PIN\_CONNECT\_USE & LIBRARIES\_AND\_VIEWS & LINE \\
LINE\_END\_DIFF\_COLOR & LINE\_END\_SAME\_COLOR & LINE\_END\_TO\_LINE\_END \\
LINE\_END\_TO\_SIDE & LINE\_TO\_LINE & LINEAR\_CONTEXT \\
LINK & LOCAL & LOGIC\_0\_NET \\
LOGIC\_1\_NET & LOOP & LOW \\
LOWER & LOWER\_CELLS & LOWEST \\
LVS & LVS\_NETLIST\_FLOW\_ICV & LVS\_NETLIST\_FLOW\_SPICE \\
LVS\_USER\_UNIT\_METER & LVS\_USER\_UNIT\_MICRON & M1\_ROUTE\_GUIDE \\
M10\_ROUTE\_GUIDE & M11\_ROUTE\_GUIDE & M12\_ROUTE\_GUIDE \\
M13\_ROUTE\_GUIDE & M14\_ROUTE\_GUIDE & M15\_ROUTE\_GUIDE \\
M2\_ROUTE\_GUIDE & M3\_ROUTE\_GUIDE & M4\_ROUTE\_GUIDE \\
M5\_ROUTE\_GUIDE & M6\_ROUTE\_GUIDE & M7\_ROUTE\_GUIDE \\
M8\_ROUTE\_GUIDE & M9\_ROUTE\_GUIDE & MACRO\_CELL \\
MACRO\_DESIGN & MACRO\_PIN\_CONNECT & MACRO\_PIN\_CONNECT\_USE \\
MAIN & MARKER\_LAYER\_TYPE & MASK\_EIGHT \\
MASK\_ELEVEN & MASK\_FIFTEEN & MASK\_FIVE \\
MASK\_FOUR & MASK\_FOURTEEN & MASK\_NINE \\
MASK\_ONE & MASK\_ONE\_HARD & MASK\_ONE\_SOFT \\
MASK\_SEVEN & MASK\_SIX & MASK\_TEN \\
MASK\_THIRTEEN & MASK\_THREE & MASK\_TWELVE \\
MASK\_TWO & MASK\_TWO\_HARD & MASK\_TWO\_SOFT \\
MAX & MAX\_AREA & MAX\_CONDITION \\
MAX\_CONDITION & MAX\_INSERTION & MEEF \\
MEGA & MERGE & MERGED \\
MERGE\_ALL & MERGE\_CONNECTED & MERGE\_CONNECTED\_AND\_TOP \\
MERGE\_TOP\_THEN\_CONNECTED & MERGE\_TOP\_PORTS\_THEN\_CONNECTED & MERGED \\
METAL\_HEIGHT & METER & MICRO \\
MICRON & MID\_CELL & MILKYWAY \\
MILLI & MIN & MIN\_AREA \\
MISSING\_CELLS & MISSING\_COLOR & MIXED \\
MODULE\_CELL & MODULE\_DESIGN & MORE\_THAN\_TRIPLE\_HEIGHT\_CELL \\
MUTUAL & MUTUAL\_NON\_ORTHOGONAL & NAME\_MATCHED \\
NANO & NDB\_ALL & NDB\_ANY \\
NDB\_BJT & NDB\_CACHE & NDB\_CAP \\
NDB\_CELL & NDB\_DIODE & NDB\_IGNORE \\
NDB\_IND & NDB\_INST & NDB\_MOS \\
NDB\_NET & NDB\_NETLIST & NDB\_NMOS \\
NDB\_NONE & NDB\_NP & NDB\_NPN \\
NDB\_ONE\_WAY & NDB\_PIN & NDB\_PMOS \\
NDB\_PN & NDB\_PNP & NDB\_RES \\
NDB\_SUBCKT & NDB\_UNDEFINED & NDB\_TWO\_WAY \\
NDB\_UNCACHED & NDM & NEGATIVE\_45 \\
NEIGHBORING\_GAP & NEITHER & NET\_NAME \\
NET\_PORTPROPERTY & NETLIST\_ANNOTATION\_FILE\_SPICE & NETLIST\_ONLY \\
NETLIST\_PEX\_SPICE & NETLIST\_XTR\_SPICE & NETS \\
NEVER & NEW\_AREF\_ONLY & NINETY \\
NMOS & NO\_COLOR & NO\_EXPLODE \\
NODE & NOMINAL & NOMINAL\_VARIATION \\
NON\_CONT\_PG & NON\_ORTHOGONAL & NON\_RECTANGULAR\_SHAPES \\
NONE & NONE\_DENSITY & NONE\_DENSITY\_PERIMETER \\
NONE\_INCLUSIVE & NORMAL & NOT\_A\_DOUBLE \\
NOT\_ADJACENT & NOT\_CONTAINED & NOT\_CORNER \\
NOT\_DEFINED\_LAYER\_TYPE & NOT\_TEXTED & NOTCH \\
NP & NPN & NULL\_CONNECT\_DATABASE \\
NWELL\_NET & NULL\_PARAMETRIX\_LAYER & OASIS \\
OCTAGONAL & ODD\_LOOP\_ONE\_LINK & ODD\_LOOP\_ALL\_LINKS \\
ODD\_LOOP\_ALL\_NODES & ODD\_LOOP\_ALL\_LOOP & ODD\_LOOP\_INSIDE\_RING \\
OFF & OLD\_AREF\_ONLY & ON \\
ONE & ONE\_EIGHTY & ONE\_OR\_NEITHER \\
ONE\_POLYGON\_PER\_LAYER & ONE\_X & ONLY\_COLOR\_1 \\
ONLY\_COLOR\_2 & ONLY\_COORDINATES & OPC \\
OPC\_USE & OPEN\_OUTER\_BOUNDARY & OPENACCESS \\
ORTHOGONAL & OTHER\_NETS & OUT \\
OUTPUT\_CELL & OUTPUT\_ERRORS & OUTSIDE \\
OVERLAP & OVER\_EXPOSURE & OVERSIZE \\
OVERWRITE & OWNER\_NET & OXIDE\_HEIGHT \\
PAD\_DESIGN & PAD\_FILLER\_CELL & PAD\_SPACER\_DESIGN \\
PARALLEL & PARALLEL\_POINT\_PROJECTION & PARTITION \\
PARTITION\_CONTEXT & PATH & PATH\_OBJECT \\
PATH\_SEGMENT & PAUSE & PERCENT \\
PERIMETER & PERIMETER\_DENSITY & PERPENDICULAR \\
PG\_FOLLOW\_PIN & PG\_PIN & PG\_PIN\_ONLY\_CELL \\
PG\_RING & PG\_STRIPE & PHYSICAL\_ONLY\_DESIGN \\
PICO & PIN & PIN\_BLOCKAGE \\
PIN\_NAME & PIN\_OBJECT & PIN\_TEXT \\
PINS & PLACEMENT\_BLOCKAGE & PLACEMENT\_HARD\_BLOCKAGE \\
PLACEMENT\_HARD\_MACRO\_BLOCKAGE & PLACEMENT\_PARTIAL\_BLOCKAGE & PLACEMENT\_SOFT\_BLOCKAGE \\
PM\_EXACT & PM\_EDGE\_NONUNIFORM & PM\_EDGE\_UNIFORM \\
PM\_EQ & PM\_FUZZY & PM\_GE \\
PM\_GT & PM\_LE & PM\_LT \\
PM\_MARKER\_CENTER & PM\_MARKER\_EDGE & PM\_NONE \\
PM\_ONE\_DIMENSIONAL & PM\_TWO\_DIMENSIONAL & PM\_TWO\_DIMENSIONAL\_RUN\_LENGTH \\
PM\_UNIFORM & PMOS & PN \\
PNP & POINT & POINT\_TO\_POINT \\
POINT\_TOUCH & POLY\_ROUTE\_GUIDE & POLYGON \\
POLYGON\_OBJECT & POLYGON\_TEXT & PORT\_NAME \\
PORT\_TEXT\_MATCHED & PORT\_TOPOLOGY\_MATCHED & POSITIVE\_45 \\
POWER & POWER\_AND\_GROUND & POWER\_NET \\
POWER\_OR\_GROUND & POWER\_XOR\_GROUND & PRE\_COLOR\_ALL\_NODES \\
PRE\_COLOR\_ENVELOPE\_PATH & PRE\_COLOR\_MIN\_PATH & PRE\_THREE\_COLOR\_ALL\_NODES \\
PRE\_THREE\_COLOR\_MIN\_PATH & PREFIX\_NAME & PRIMARY\_AXIS \\
PRINT & PRINT\_ALL & PRINT\_RESET\_ALL \\
PROJECTING & PROJECTION & PROPERTY \\
PROPERTY\_ERRORS\_ONLY & PROPERTY\_TEXT & PWELL\_NET \\
QCMP & QTF\_LAYER\_TYPE & RADIAL \\
RADIAL\_INSIDE & RADIAL\_INTERSECTION & RADIAL\_OUTSIDE \\
RAM\_CELL & REASSIGN & REASSIGNED\_SHORTED \\
RDL\_USE & RECT & RECTANGLE \\
RECTANGLE\_OBJECT & REDUCTION & REGION \\
REGION\_CONTEXT & REGULAR\_TEXT & RELATED\_COINCIDENT \\
RELATIVE & REMAINDER & REMOVE\_LAYER\_TYPE \\
REMOVE\_MULTI\_EQUIVS & RENAME & RENAMED \\
REPLACE & REPORT & RESET \\
RESET\_NET & RESISTOR & RETAIN\_45 \\
RIGHT & RING & RING\_USE \\
ROM\_CELL & ROTATE\_180 & ROTATIONAL \\
ROUND & ROUTE\_GUIDE\_OBJECT & ROUTE\_TYPE \\
ROUTING\_BLOCKAGE & ROUTING\_FOR\_TOP\_LEVEL\_BLOCKAGE & RPGROUP\_BLOCKAGE \\
SAME & SAME\_COLOR & SAME\_DEVICE\_NAME\_ONLY \\
SAME\_MASK & SAME\_NET & SAME\_POLYGON \\
SCAN\_NET & SCREEN\_BLOCKAGE & SERIES \\
SERIES\_PARALLEL & SHAPING\_BLOCKAGE & SHIELD\_ROUTE \\
SHIELD\_ROUTE\_USE & SHORT\_BY\_LONG & SHORTED \\
SHRINK & SIDE\_DIFF\_COLOR & SIDE\_SAME\_COLOR \\
SIDE\_TO\_SIDE & SIGNAL & SIGNAL\_CONTEXT \\
SIGNAL\_DETAIL & SIGNAL\_GLOBAL & SIGNAL\_LOCAL \\
SIGNAL\_NET & SIGNAL\_USER & SIMPLE \\
SIMULATE & SINGLE & SINGLE\_DEVICE\_EQUIVS \\
SINKS & SIXTEEN & SKIP \\
SKIP\_PARENTS & SLASH & SLEEP\_CONTROL \\
SLOT\_BLOCKAGE & SOFT\_MACRO\_IO\_FIXED\_CELL & SOG\_DESIGN\_CELL \\
SOURCES\_ON\_SAME\_NET & SPECIAL\_POWER & SPECIAL\_VIA\_CELL \\
SPICE & SPICE\_BJT & SPICE\_CAP \\
SPICE\_CELL & SPICE\_DIODE & SPICE\_IND \\
SPICE\_MOS & SPICE\_RES & SPLIT\_CLOCK \\
SQUARE & SQUARE\_CENTERED & SQUARE\_EXTENDED \\
SQUARE\_INSIDE & SQUARE\_OUTSIDE & STACK\_VIA\_CELL \\
STANDARD\_CELL & STANDARD\_FILLER\_CELL & START \\
STATISTICAL\_COUNT & STOP & STRING \\
STRIPE & STRIPE\_USE & SUBSTRING\_NAME \\
SUM & SYMMETRIC & SYMMETRIC\_NON\_INTERSECTING \\
TAP\_CELL & TECH & TERMINAL \\
TERMINAL\_NAME & TERMINAL\_TEXT & TEXT \\
TEXT\_DISPLAY & TEXTED & TIE\_HIGH \\
TIE\_LOW & TOP & TOP\_CELL \\
TOP\_OF\_NET & TOP\_PORT & TOUCH \\
TREE & TRIANGLE & TRIANGLE\_ONE\_LINK \\
TRIANGLE\_ALL\_LINKS & TRIANGLE\_ALL\_NODES & TRIPLE\_HEIGHT\_STANDARD\_CELL \\
TRUNCATE & TSV\_CELL & TWO \\
TWO\_SEVENTY & TWO\_X & UF\_ADJUSTABLE \\
UF\_CELL & UF\_CHECKERED & UF\_DPT\_SPACING \\
UF\_EXPANDABLE & UF\_LAYER & UF\_LINE\_END \\
UF\_LINE\_SIDE & UF\_PATTERN & UF\_POLYGON \\
UF\_STACK & UF\_STRIPE & UNCONNECTED \\
UNDER\_EXPOSURE & UNDERSIZE & UNSET\_NET \\
UNKNOWN\_L\_MODEL\_CELL & UNMATCHED\_PINS & UNMATCHED\_TEXT \\
UNSPECIFIED & UNSPECIFIED\_DOUBLE & UNUSED \\
UNUSED\_CELLS & UNUSED\_MARKERS & UPDATE \\
UPPER & USE & USE\_ALL\_EQUIV \\
USE\_MILKYWAY & USE\_OPENACCESS & USE\_NDM \\
USED\_CELLS & USER\_DEFINED & USER\_ROUTE \\
USER\_ROUTE\_USE & VCELL & VCMP \\
VERBOSE & VERILOG & VERTICAL \\
VERTICAL\_AXIS & VERTICAL\_TRACK & VERTICAL\_TRACK\_OBJECT \\
VERTICAL\_WIRE\_TRACK & VERTICALWIRE & VERTICES \\
VIA\_BLOCKAGE & VIA\_LAYER\_TYPE & VIA\_OBJECT \\
VIEWONLY\_LAYER\_TYPE & VIEWS & VIRTUAL\_AMBIT \\
VIRTUAL\_CORE & WAIVE & WARN \\
WARNING & WATCH & WELL\_TAP\_DESIGN \\
WINDOW & WIRING\_BLOCKAGE & WITH\_PORTS \\
X\_BY\_Y & XO\_CELL & ZERO\_SKEW \\
ZERO\_SKEW\_USE &  &  \\
\end{longtable}

\subsection{运算符}
\subsubsection*{Operators}

运算符是具有特定含义的符号。例如，\cmd{+} 表示将两个数相加。运算符类型见表~\ref{tab:pxl-op-types}。

\begin{longtable}{@{}p{0.22\textwidth} p{0.28\textwidth} p{0.42\textwidth}@{}}
\caption{运算符类型（Operator Types）}\label{tab:pxl-op-types}\\
\toprule
\textbf{运算符类型} & \textbf{操作数个数} & \textbf{示例} \\
\midrule
\endfirsthead
\toprule
\textbf{运算符类型} & \textbf{操作数个数} & \textbf{示例} \\
\midrule
\endhead
\bottomrule
\endfoot
一元（unary） & 1 & 逻辑非（\cmd{!}） \\
二元（binary） & 2 & 加法（\cmd{+}）、幂运算（\cmd{**}） \\
三元（ternary） & 3 & 条件运算符（\cmd{e1 ? e2 : e3}） \\
\end{longtable}

\subsection{注释}
\subsubsection*{Comments}

注释是用于说明程序的文本，便于向读者表明代码意图。注释按空白处理。表~\ref{tab:pxl-comment-types} 按注释所占行数给出注释类型。

\begin{longtable}{@{}p{0.22\textwidth} >{\raggedright\arraybackslash\ttfamily}p{0.18\textwidth} >{\raggedright\arraybackslash\ttfamily}p{0.22\textwidth} p{0.28\textwidth}@{}}
\caption{注释类型（Comment Types）}\label{tab:pxl-comment-types}\\
\toprule
\textbf{\textrm{注释类型}} & \textbf{\textrm{起始}} & \textbf{\textrm{结束}} & \textbf{\textrm{长度}} \\
\midrule
\endfirsthead
\toprule
\textbf{\textrm{注释类型}} & \textbf{\textrm{起始}} & \textbf{\textrm{结束}} & \textbf{\textrm{长度}} \\
\midrule
\endhead
\bottomrule
\endfoot
\textrm{行注释（line）} & // & 下一换行 & 一行 \\
\textrm{块注释（block）} & /* & */ & 任意长度 \\
\end{longtable}

块注释不可嵌套。若在注释内看到 \cmd{*/}，预处理器会报告警告。例如，若 runset 含有：
\begin{lstlisting}[style=pxl]
/*           - 1st comment start
       /*       - 2nd comment start
       */       - 2nd comment end
*/           - 1st comment end
\end{lstlisting}

则「1st comment start」与「2nd comment end」配对；「1st comment start」不会与「1st comment end」配对。

注释示例如下：
\begin{lstlisting}[style=pxl]
/* this is a single-line block comment */

// this is a line comment, so /* is ignored

/* block comments /* cannot nest. Therefore this comment evokes a
warning from the preprocessor */

/*
 * block comments can span
 * several lines, such
 * as this one does.
 */

/*
  A leading * is not required on each line of a
  block comment but does make the comment stand out
  from the program.
*/
\end{lstlisting}

\subsection{环境变量引用}
\subsubsection*{Environment Variable Reference}

环境变量是一组可用于控制进程运行方式的动态值。

环境变量引用由美元符（\cmd{\$}）后跟任意合法标识符组成。若引用未定义的环境变量，将在编译期报错并拒绝。

环境变量会被替换为包含其值的字符串字面量。仅对形如 \cmd{\$<id>} 的记号进行替换。具体而言，字符串字面量内部不会替换；由于相邻字符串字面量会拼接，通常也无需在字面量内替换。

\begin{noteBox}
若环境变量的值为空字符串，则将其转换为值为空字符串的字符串字面量。
\end{noteBox}

环境变量用法示例如下：
\begin{lstlisting}[style=pxl]
x : string = $y;

/*
 * assuming environment variable my_value contains the value abcd,
 * the following code evaluates to true; otherwise evaluates to false.
 */
res : boolean = ( $my_value == "abcd" );

/*
 * if environment variable aaa contains the value
 * /home/dir/path, the following code produces a string
 * containing a path to the file "cases" in that directory.
 */
path : string = $aaa "/cases";

/*
 * if environment variable bbb is defined but empty, the
 * following code produces the string abcdef.
 */
mystr : string = "abc" $bbb "def";
\end{lstlisting}

% ============================================================
\section{变量与运算符}
\subsection*{Variables and Operators}

本节介绍 PXL 中的变量与运算符概念。

\subsection{变量}
\subsubsection*{Variables}

变量用于在程序中临时存储信息。
声明变量时，变量名后须跟冒号（\cmd{:}）与变量类型。声明时也可给出初值。见下一小节示例。

\begin{noteBox}
有关命名限制，见「标识符」小节。
\end{noteBox}

\subsubsection{初始化变量}
\paragraph*{Initializing Variables}

可用两种方式将变量初始化为任意受支持的数据类型：
\begin{itemize}
  \item 使用字面量值。
\begin{lstlisting}[style=pxl]
y : boolean = true;                   // initialize by literal value
\end{lstlisting}
  \item 使用表达式。
\begin{lstlisting}[style=pxl]
z : boolean = !y;                     // initialized by expression
\end{lstlisting}
\end{itemize}

下列示例展示不同的初始化方法：
\begin{lstlisting}[style=pxl]
x : integer;                              // no initial value
pi : double = 3.141592;                   // initialized by value
pi2 : double = pi * 2;                    // initialized by expression
s : string = "abc" "def";                 // concatenated
s1 : string = "abc" + pi;                 // converted to string,then concatenated
c : constraint of double = (10.0,200];    // initialized by value
\end{lstlisting}

\subsubsection{变量作用域}
\paragraph*{Scope of Variables}

PXL 中的作用域（scope）指变量与表达式在程序若干上下文中的含义。下文通过示例说明与变量作用域相关的规则与条件。

局部作用域由一对花括号（\cmd{\{\}}）定义。也就是说，用花括号定义的复合语句会引入新的变量作用域。花括号的用法见后文「复合语句」（Compound Statement）一节。

流程控制语句会创建新作用域。
\begin{lstlisting}[style=pxl]
if ( var == 1) {                // Start of new scope
   my_var: integer = 1;
   ...
}                             // End of scope.
\end{lstlisting}

\begin{itemize}
  \item 在局部作用域中声明的变量在该作用域外不存在。也就是说，若变量在某一作用域内声明，则不能在该作用域外使用。例如：
\begin{lstlisting}[style=pxl]
{
   x = y + b;
}
z = x;      // x is not defined; it only exists in the local scope
\end{lstlisting}
  \item 变量从最外层作用域到最内层作用域全局可见。例如：
\begin{lstlisting}[style=pxl]
{
    x : integer = 1;
       {
       a = x + b; // x exists and can be seen within this scope
                   // (global from outer scope)
       x = y + b; // global value of x is overwritten
       }
    z = x;         // x is defined here and contains the value of y + b
}
z = x;                // x is not defined here - parser error
\end{lstlisting}
  \item 在外层作用域声明、又在内层作用域重新声明的变量会被遮蔽（shadowed），以避免值冲突。
  \begin{itemize}
    \item 若发生变量遮蔽，且被遮蔽变量的作用域并非由函数声明定义，则程序生成警告。
    \item 在函数调用内，遮蔽是预期行为，不会警告。例如：
\begin{lstlisting}[style=pxl]
y: integer = 1;
x: integer = 1;
{
   y : integer = 3;      // The original y value is shadowed.
                         // This is a new y.
    x = x + 1;           // x is equal to 2 here
    y = y + 1;           // y is equal to 4 here
}
                         // y is equal to 1 here - the shadowed value
                         // returns when exiting the inner scope
                         // x is equal to 2 here
\end{lstlisting}
  \end{itemize}
\end{itemize}

下面是变量遮蔽的另一示例。函数声明：
\begin{lstlisting}[style=pxl]
fxy : function (
   x : integer;
   y : integer;
) returning z : integer {
      x = x + 1;
      y = y + 1;
      z = x + y;
};
\end{lstlisting}

函数调用：
\begin{lstlisting}[style=pxl]
x1 : integer = 1;
y : integer = 1;

z1 = fxy(x1, y);                              // within the function y is shadowed
                                              // y in the local scope is "2"
x2 = x1;                                      // x1 is "1"
y1 = y;                                       // y is still "1"
z2 = z1;                                      // z1 is "4"
\end{lstlisting}

\subsection{运算符}
\subsubsection*{Operators}

运算符是表达式中用于指定在求值时应执行某些操作的符号。

PXL 支持下列运算符：
\begin{itemize}
  \item 赋值运算符（Assignment Operators）
  \item 关系运算符（Relational Operators）
  \item 逻辑运算符（Logical Operators）
  \item 条件运算符（Conditional Operator）
  \item 数值与按位运算符（Numeric and Bitwise Operators）
  \item 仅输出 double 的运算符（Operators With Double-Only Outputs）
\end{itemize}

\subsubsection{赋值运算符}
\paragraph*{Assignment Operators}

赋值运算符对表达式求值并将其保存到目标左值（lvalue）。
左值是可容纳值、并可出现在赋值语句左侧的表达式。

涉及在约束值中使用 \cmd{>} 或 \cmd{>=} 的赋值，赋值运算符前后必须有空白。也可将约束值放在赋值 \cmd{=} 之后的括号内。

例如：
\begin{lstlisting}[style=pxl]
abc = >=3;
xyz = (>5);
\end{lstlisting}

\begin{noteBox}
赋值运算符不可串联。下例不合法：\\
\cmd{a = b = c}
\end{noteBox}

正确赋值示例如下：
\begin{lstlisting}[style=pxl]
x:      integer;      // declare a variable named x

point: newtype struct of { x: integer; y: integer; }; // declare a struct
                                                      // of integers
p: point;
x   = 42;      // assign x a value of 42
p.x = 10;      // assign p.x a value of 10
p.y = 27;      // assign p.y a value of 27
\end{lstlisting}

\subsubsection{关系运算符}
\paragraph*{Relational Operators}

关系表达式将表达式求值为布尔值 \cmd{true} 或 \cmd{false}。表~\ref{tab:pxl-rel-ops} 定义关系运算符。

\begin{longtable}{@{}>{\raggedright\arraybackslash\ttfamily}p{0.16\textwidth} p{0.42\textwidth} p{0.34\textwidth}@{}}
\caption{关系运算符（Relational Operators）}\label{tab:pxl-rel-ops}\\
\toprule
\textbf{\textrm{语法}} & \textbf{\textrm{说明}} & \textbf{\textrm{适用于}} \\
\midrule
\endfirsthead
\toprule
\textbf{\textrm{语法}} & \textbf{\textrm{说明}} & \textbf{\textrm{适用于}} \\
\midrule
\endhead
\bottomrule
\endfoot
e1 == e2 & 若 e1 等于 e2 则为 true，否则为 false & 所有完全定义的数据类型 \\
e1 != e2 & 若 e1 不等于 e2 则为 true，否则为 false & 所有完全定义的数据类型 \\
e1 < e2 & 若 e1 小于 e2 则为 true，否则为 false & 数值 \\
e1 <= e2 & 若 e1 不大于 e2 则为 true，否则为 false & 数值 \\
e1 >= e2 & 若 e1 不小于 e2 则为 true，否则为 false & 数值 \\
e1 > e2 & 若 e1 大于 e2 则为 true，否则为 false & 数值 \\
\end{longtable}

\subsubsection{逻辑运算符}
\paragraph*{Logical Operators}

逻辑表达式求值为布尔值。逻辑运算符要求全部参数均为 \cmd{boolean} 类型。表达式从左到右求值，并在结果已确定时停止。表~\ref{tab:pxl-logic-ops} 定义逻辑运算符。

\begin{longtable}{@{}>{\raggedright\arraybackslash\ttfamily}p{0.22\textwidth} p{0.70\textwidth}@{}}
\caption{逻辑运算符（Logical Operators）}\label{tab:pxl-logic-ops}\\
\toprule
\textbf{\textrm{语法}} & \textbf{\textrm{说明}} \\
\midrule
\endfirsthead
\toprule
\textbf{\textrm{语法}} & \textbf{\textrm{说明}} \\
\midrule
\endhead
\bottomrule
\endfoot
e1 \&\& e2 & 若 e1 为 false，则结果为 false，否则为 e2 \\
e1 || e2 & 若 e1 为 true，则结果为 true，否则为 e2 \\
! e1 & 若 e1 为 true，则结果为 false，否则为 true \\
\end{longtable}

\subsubsection{条件运算符}
\paragraph*{Conditional Operator}

条件表达式接受三个参数；第一个参数必须为 \cmd{boolean} 类型。表~\ref{tab:pxl-cond-op} 定义条件运算符。

\begin{longtable}{@{}>{\raggedright\arraybackslash\ttfamily}p{0.22\textwidth} p{0.70\textwidth}@{}}
\caption{条件运算符（Conditional Operator）}\label{tab:pxl-cond-op}\\
\toprule
\textbf{\textrm{语法}} & \textbf{\textrm{说明}} \\
\midrule
\endfirsthead
\toprule
\textbf{\textrm{语法}} & \textbf{\textrm{说明}} \\
\midrule
\endhead
\bottomrule
\endfoot
e1 ? e2 : e3 & 若 e1 为 true，则为 e2，否则为 e3 \\
\end{longtable}

考虑带有三个参数 \cmd{e1}、\cmd{e2}、\cmd{e3} 的条件运算符。求值顺序如下：
\begin{enumerate}
  \item 求值第一个参数（\cmd{e1}）。
  \item 若第一个参数（\cmd{e1}）为 \cmd{true}，则求值第二个参数。
  \item 若第一个参数（\cmd{e1}）为 \cmd{false}，则求值第三个参数（\cmd{e3}）。
\end{enumerate}

\subsubsection{数值与按位运算符}
\paragraph*{Numeric and Bitwise Operators}

PXL 支持全部标准数学运算符。按位运算符仅支持整数。表~\ref{tab:pxl-num-ops} 定义数值与按位运算符。

\begin{longtable}{@{}>{\raggedright\arraybackslash\ttfamily}p{0.14\textwidth} p{0.22\textwidth} p{0.24\textwidth} p{0.32\textwidth}@{}}
\caption{数值与按位运算符（Numeric and Bitwise Operators）}\label{tab:pxl-num-ops}\\
\toprule
\textbf{\textrm{语法}} & \textbf{\textrm{说明}} & \textbf{\textrm{允许类型}} & \textbf{\textrm{注释}} \\
\midrule
\endfirsthead
\toprule
\textbf{\textrm{语法}} & \textbf{\textrm{说明}} & \textbf{\textrm{允许类型}} & \textbf{\textrm{注释}} \\
\midrule
\endhead
\bottomrule
\endfoot
e1 + e2 & 加法 & Integer 与 double & 任一参数为 double 时返回 double \\
e1 - e2 & 减法 & Integer 与 double & 任一参数为 double 时返回 double \\
e1 * e2 & 乘法 & Integer 与 double & 任一参数为 double 时返回 double \\
e1 \% e2 & 整数余数 & Integer 与 double & 任一参数为 double 时返回 double \\
e1 << e2 & 逻辑左移（移入零） & Integer &  \\
e1 >> e2 & 逻辑右移（移入零） & Integer &  \\
e1 \& e2 & 按位 AND & Integer &  \\
e1 | e2 & 按位 OR & Integer &  \\
e1 \^{} e2 & 按位 XOR & Integer &  \\
! e1 & 按位 NOT & Integer &  \\
\end{longtable}

\subsubsection{仅输出 double 的运算符}
\paragraph*{Operators With Double-Only Outputs}

除法（\cmd{/}）与幂运算（\cmd{**}）将参数视为 double，并且始终只返回 double。表~\ref{tab:pxl-double-ops} 定义这些运算符。

\begin{longtable}{@{}>{\raggedright\arraybackslash\ttfamily}p{0.14\textwidth} p{0.20\textwidth} p{0.24\textwidth} p{0.34\textwidth}@{}}
\caption{仅输出 double 的运算符（Operators With Double-Only Outputs）}\label{tab:pxl-double-ops}\\
\toprule
\textbf{\textrm{语法}} & \textbf{\textrm{说明}} & \textbf{\textrm{接受}} & \textbf{\textrm{注释}} \\
\midrule
\endfirsthead
\toprule
\textbf{\textrm{语法}} & \textbf{\textrm{说明}} & \textbf{\textrm{接受}} & \textbf{\textrm{注释}} \\
\midrule
\endhead
\bottomrule
\endfoot
e1 / e2 & 除法 & Integer 与 double & e1 与 e2 被强制为 double；运算始终返回 double \\
e1 ** e2 & 幂运算 & Double & e1 与 e2 被强制为 double；运算始终返回 double \\
\end{longtable}

\begin{noteBox}
任何含除法运算符（\cmd{/}）或幂运算符（\cmd{**}）的表达式，均不可用作列表或哈希表的索引。
\end{noteBox}

若要执行整数运算，必须将该运算包在 \cmd{dtoi()} 函数调用中。该函数将 double 转换为 integer。例如：
\begin{lstlisting}[style=pxl]
var : integer = dtoi(10/3);
\end{lstlisting}

\subsection{成员引用}
\subsubsection*{Member Reference}

成员引用表达式访问复合对象的一个成员。成员引用取该单一成员的类型。该成员是左值，并可作为赋值目标。表~\ref{tab:pxl-member-ref} 定义成员引用表达式。
有关左值的更多信息，见「赋值运算符」。

\begin{longtable}{@{}>{\raggedright\arraybackslash\ttfamily}p{0.22\textwidth} p{0.70\textwidth}@{}}
\caption{成员引用（Member Reference）}\label{tab:pxl-member-ref}\\
\toprule
\textbf{\textrm{语法}} & \textbf{\textrm{说明}} \\
\midrule
\endfirsthead
\toprule
\textbf{\textrm{语法}} & \textbf{\textrm{说明}} \\
\midrule
\endhead
\bottomrule
\endfoot
e1 . e2 & 按名称访问结构体成员 \\
e1 [ e2 ] & 按索引访问列表或哈希成员 \\
\end{longtable}

\subsection{优先级与结合性}
\subsubsection*{Precedence and Associativity}

表~\ref{tab:pxl-precedence} 列出 PXL 中全部运算符的优先级与结合性。表顶优先级最高，表底最低。

\begin{longtable}{@{}>{\raggedright\arraybackslash\ttfamily}p{0.48\textwidth} p{0.44\textwidth}@{}}
\caption{运算符的优先级与结合性（最高优先级在前）}\label{tab:pxl-precedence}\\
\toprule
\textbf{\textrm{运算符}} & \textbf{\textrm{结合性}} \\
\midrule
\endfirsthead
\toprule
\textbf{\textrm{运算符}} & \textbf{\textrm{结合性}} \\
\midrule
\endhead
\bottomrule
\endfoot
() \quad [] \quad . & 从左到右 \\
** & 从右到左（幂运算） \\
! \quad + \quad - & 从右到左（一元运算符） \\
== \quad != \quad < \quad <= \quad > \quad >= & 非法（一元运算符） \\
* \quad / \quad \% & 从左到右 \\
+ \quad - & 从左到右 \\
<< \quad >> & 从左到右 \\
< \quad <= \quad > \quad >= & 非结合 \\
== \quad != & 非结合 \\
\& & 从左到右 \\
\^{} & 从左到右 \\
| & 从左到右 \\
\&\& & 从左到右 \\
|| & 从左到右 \\
?: & 从右到左 \\
= \quad @= \quad ?= & 非法 \\
\textrm{用户定义运算符} & 非结合 \\
\end{longtable}

表~\ref{tab:pxl-associativity} 定义各类结合性。

\begin{longtable}{@{}p{0.28\textwidth} p{0.64\textwidth}@{}}
\caption{结合性（Associativity）}\label{tab:pxl-associativity}\\
\toprule
\textbf{结合性} & \textbf{定义} \\
\midrule
\endfirsthead
\toprule
\textbf{结合性} & \textbf{定义} \\
\midrule
\endhead
\bottomrule
\endfoot
从左到右（Left to right） & 表达式从左到右求值。 \\
从右到左（Right to left） & 表达式从右到左求值。 \\
非结合（Nonassociative） & 需要括号。 \\
非法（Illegal） & 该运算符不可串联。运算符产生的类型与其消费的类型不同。 \\
\end{longtable}

下列示例说明结合性如何工作：
\begin{lstlisting}[style=pxl]
/*
 * pairs of expressions show evaluation order
 */

A + B + C
( A + B ) + C

A ** B ** C
A ** ( B ** C )

A * B + C >> 2
( ( A * B ) + C ) >> 2

A >= B == C
( A >= B ) == C

A < B == C > D
( A < B ) == ( C > D )
\end{lstlisting}

% ============================================================
\section{数据类型与类型系统}
\subsection*{Data Types and Typing}

本节说明 PXL 支持的数据类型，以及决定 PXL 内类型如何工作的因素。

变量可显式声明类型。例如：
\begin{lstlisting}[style=pxl]
i : integer = 1 ;           // i is declared as an integer
\end{lstlisting}

变量也可根据运算符的返回类型隐式定型。例如：
\begin{lstlisting}[style=pxl]
a : double = 1.1203;
b : double = 2.0213;
c = a + b;                     // c is implicitly the type of "double"
\end{lstlisting}

\subsection{受支持的数据类型}
\subsubsection*{Supported Data Types}

PXL 支持多数基本类型以及一些工具相关数据类型。

\subsubsection{基本类型}
\paragraph*{Primitive Types}

表~\ref{tab:pxl-primitive} 给出 PXL 支持的基本数据类型。

\begin{longtable}{@{}>{\raggedright\arraybackslash\ttfamily}p{0.34\textwidth} p{0.58\textwidth}@{}}
\caption{基本类型（Primitive Types）}\label{tab:pxl-primitive}\\
\toprule
\textbf{\textrm{数据类型}} & \textbf{\textrm{说明}} \\
\midrule
\endfirsthead
\toprule
\textbf{\textrm{数据类型}} & \textbf{\textrm{说明}} \\
\midrule
\endhead
\bottomrule
\endfoot
integer & 整数；大小由实现定义，一般为 32 位有符号值。 \\
double & 双精度浮点；大小在实现时定义，一般为 64 位值。 \\
boolean & true 或 false \\
string & 字符序列；支持用转义序列包含特殊字符。 \\
handle & 不透明数据类型，用于跟踪 PXL 程序本身之外的数据对象。 \\
constraint of double 与 constraint of integer & 连续 double 值或 integer 值集合的简洁表示。 \\
predicate\_handle & 布尔表达式，并对将层标识符求值为空/非空有特殊处理。 \\
\end{longtable}

\paragraph{integer}

整数声明示例如下：
\begin{lstlisting}[style=pxl]
i : integer;     // declared but not initialized
i : integer = 1; // declared and initialized
\end{lstlisting}

\paragraph{string}

定义零个或多个字符的字符串。\cmd{string} 类型由字符串字面量初始化。例如：
\begin{lstlisting}[style=pxl]
s : string = "xyzzy";
\end{lstlisting}

\cmd{string} 类型对任意 \cmd{string} 类型表达式支持表~\ref{tab:pxl-string-ops} 所示操作。

\begin{longtable}{@{}>{\raggedright\arraybackslash\ttfamily}p{0.18\textwidth} p{0.12\textwidth} p{0.12\textwidth} p{0.14\textwidth} p{0.34\textwidth}@{}}
\caption{字符串类型操作（String Type Operations）}\label{tab:pxl-string-ops}\\
\toprule
\textbf{\textrm{操作}} & \textbf{\textrm{e2 类型}} & \textbf{\textrm{返回}} & \textbf{\textrm{修改 e1}} & \textbf{\textrm{说明}} \\
\midrule
\endfirsthead
\toprule
\textbf{\textrm{操作}} & \textbf{\textrm{e2 类型}} & \textbf{\textrm{返回}} & \textbf{\textrm{修改 e1}} & \textbf{\textrm{说明}} \\
\midrule
\endhead
\bottomrule
\endfoot
e1.empty() & -- & boolean & 否 & 若字符串为空（即 \cmd{""}）则返回 true；否则返回 false。 \\
e1.size() & -- & integer & 否 & 返回字符串中的字符数。 \\
\end{longtable}

\paragraph{handle}

外部管理对象的内部表示。类型 \cmd{handle} 具有对每个外部管理对象唯一的字符串表示。

\paragraph{constraint of double 与 constraint of integer}

\cmd{constraint of double} 与 \cmd{constraint of integer} 是尺寸规则值的数学区间。这些约束为描述 DRC 间距距离等参数提供直观机制。

若未指定约束类型，则隐含该区间为 \cmd{constraint of double}。

约束数据类型对属于该类型的表达式支持表~\ref{tab:pxl-constraint-ops} 所示操作。

\begin{longtable}{@{}>{\raggedright\arraybackslash\ttfamily}p{0.20\textwidth} p{0.12\textwidth} p{0.14\textwidth} p{0.14\textwidth} p{0.30\textwidth}@{}}
\caption{constraint of double 与 constraint of integer 的操作}\label{tab:pxl-constraint-ops}\\
\toprule
\textbf{\textrm{操作}} & \textbf{\textrm{e2 类型}} & \textbf{\textrm{返回}} & \textbf{\textrm{修改 e1}} & \textbf{\textrm{说明}} \\
\midrule
\endfirsthead
\toprule
\textbf{\textrm{操作}} & \textbf{\textrm{e2 类型}} & \textbf{\textrm{返回}} & \textbf{\textrm{修改 e1}} & \textbf{\textrm{说明}} \\
\midrule
\endhead
\bottomrule
\endfoot
e1.contains(e2) & item & boolean & 否 & 若 e2 值落在 e1 约束指定的区间内则返回 true。 \\
e1.category() & -- & enum & 否 & 返回该约束在预定义枚举 constraint\_category 中的枚举元。见表~\ref{tab:pxl-constraint-cat}。 \\
e1.lo() & -- & integer/double & 否 & 对二元约束，返回下界：constraint of double 为 double；constraint of integer 为 integer。\textsuperscript{1} \\
e1.hi() & -- & integer/double & 否 & 对二元约束，返回上界：constraint of double 为 double；constraint of integer 为 integer。\textsuperscript{1} \\
\end{longtable}

\noindent\small
\textsuperscript{1}对一元约束（如 \cmd{==} 与 \cmd{>}），\cmd{e1.lo()} 与 \cmd{e1.hi()} 均返回该单一界限。

约束数据类型只能通过从另一同类型约束赋值，或从字面常量创建，如表~\ref{tab:pxl-constraint-cat} 所示。

\begin{longtable}{@{}>{\raggedright\arraybackslash\ttfamily}p{0.14\textwidth} p{0.48\textwidth} >{\raggedright\arraybackslash\ttfamily}p{0.28\textwidth}@{}}
\caption{约束类别（Constraint Categories）}\label{tab:pxl-constraint-cat}\\
\toprule
\textbf{\textrm{字面量}} & \textbf{\textrm{说明}} & \textbf{\textrm{约束类别}} \\
\midrule
\endfirsthead
\toprule
\textbf{\textrm{字面量}} & \textbf{\textrm{说明}} & \textbf{\textrm{约束类别}} \\
\midrule
\endhead
\bottomrule
\endfoot
{[a,b]} & 从 a 到 b 的闭区间，包含两端点 & CONSTRAINT\_GELE \\
{(a,b)} & 从 a 到 b 的开区间，不含两端点 & CONSTRAINT\_GTLT \\
{[a,b)} & 从 a 到 b 的半开区间，含 a 不含 b & CONSTRAINT\_GELT \\
{(a,b]} & 从 a 到 b 的半开区间，不含 a 含 b & CONSTRAINT\_GTLE \\
{>= a} & 大于等于 a 的全部数 & CONSTRAINT\_GE \\
{> a} & 大于 a 的全部数 & CONSTRAINT\_GT \\
{<= a} & 小于等于 a 的全部数 & CONSTRAINT\_LE \\
{< a} & 小于 a 的全部数 & CONSTRAINT\_LT \\
{== a} & 仅包含数 a 的区间 & CONSTRAINT\_EQ \\
{!= a} & 包含除 a 以外全部数的区间 & CONSTRAINT\_NE \\
\end{longtable}

\cmd{lo()} 与 \cmd{hi()} 对所有约束均有效。对单端约束，两函数均定义为返回该单一端点 \cmd{a}。对双端约束，\cmd{lo()} 返回 \cmd{a}，\cmd{hi()} 返回 \cmd{b}。

\paragraph{示例}
\textit{Examples}

示例如下。
\begin{itemize}
  \item Constraint of integer
  \begin{itemize}
    \item \cmd{[2,8]} 包含值 2、3、4、5、6、7、8
    \item \cmd{(2,8)} 包含值 3、4、5、6、7
  \end{itemize}
  \item Constraint of double
  \begin{itemize}
    \item \cmd{[2,8]} 包含该区间内全部实数（含 2.0 与 8.0）
    \item \cmd{(2,8)} 包含除 2.0 与 8.0 外该区间内的全部实数
  \end{itemize}
  \item 声明示例
\begin{lstlisting}[style=pxl]
a_value : constraint of double = [2.0,8.0];
\end{lstlisting}
  \item 赋值示例
\begin{lstlisting}[style=pxl]
distance <= 2.0                 // All values less than or equal to 2.0
distance = <= 2.0               // All values less than or equal to 2.0
distance = [ 2.0, 8.0 )         // All values >= 2.0 and < 8.0
count = 5                       // Count is exactly equal to 5
count = ==5                     // Count is exactly equal to 5
\end{lstlisting}
\begin{noteBox}
诸如 \cmd{==} 与 \cmd{<=} 的符号是约束的一部分。对期望约束类型的变量，赋值运算符 \cmd{=} 可选。但若右侧是约束类型变量，则必须使用赋值运算符：
\end{noteBox}
\begin{lstlisting}[style=pxl]
my_constraint: constraint of double = < 5.0;
external1(layer1, distance = my_constraint );
external1(layer1, distance <5);                              // same as above
external1(layer1, distance = < 5);                           // same as above
\end{lstlisting}
\end{itemize}

\paragraph{predicate\_handle}

括在双方括号 \cmd{[[\ldots]]} 中的布尔表达式。允许任意包含至少一个层名或违例名的布尔逻辑。层名在谓词表达式内有特殊含义：运行时若该层有数据则求值为 \cmd{true}，若层为空则求值为 \cmd{false}。违例对象在存在违例时为 \cmd{true}，无违例时为 \cmd{false}。

在 runset 函数中，\cmd{predicate\_handle} 参数根据表达式求值控制给定函数的行为：
\begin{itemize}
  \item \cmd{true}：函数按正常方式运行。
  \item \cmd{false}：函数产生空输出层。
\end{itemize}

例如：
\begin{lstlisting}[style=pxl]
temp_layer = and(poly, diff, predicate = [[metal1]]);
\end{lstlisting}

若 \cmd{metal1} 为空，则跳过 \cmd{and()} 函数，且 \cmd{temp\_layer} 为空。

\subsection{用户定义结构体与类型}
\subsubsection*{User-Defined Structures and Types}

除基本类型外，PXL 还支持表~\ref{tab:pxl-user-types} 所示的结构体与类型。

\begin{longtable}{@{}>{\raggedright\arraybackslash\ttfamily}p{0.28\textwidth} p{0.64\textwidth}@{}}
\caption{用户定义类型（User-Defined Types）}\label{tab:pxl-user-types}\\
\toprule
\textbf{\textrm{用户定义类型}} & \textbf{\textrm{主要特征}} \\
\midrule
\endfirsthead
\toprule
\textbf{\textrm{用户定义类型}} & \textbf{\textrm{主要特征}} \\
\midrule
\endhead
\bottomrule
\endfoot
newtype... & 定义类型。 \\
enum of... & 定义枚举元素列表。 \\
list of... & 定义元素列表。 \\
hash of... & 将一种数据类型的元素列表转换为另一种数据类型。 \\
struct of... & 定义由多个变量组成的复合类型。 \\
function & 定义可传递给另一用户定义函数、并可像普通变量一样调用的用户定义函数。 \\
\end{longtable}

\subsubsection{newtype...}
\paragraph*{newtype...}

\cmd{newtype} 声明一种新变量类型。它并不定义变量。每种类型都是唯一的。

合法类型定义示例如下：
\begin{lstlisting}[style=pxl]
/*
 * define a new type, direction.
 */

direction: newtype enum of {
   up, down, left, right, all // trailing comma is allowed
};

/* define two distinct types */
T1: newtype integer;
T2: newtype integer;

/* define objects of the two types */
o1: T1;
o2: T2;
\end{lstlisting}

非法类型定义示例如下：
\begin{lstlisting}[style=pxl]
/* given the definitions above, o1 and o2 are of
 * different types, therefore they cannot be compared
 * and this is an error
 */

if (o1 == o2) { ... }
\end{lstlisting}

下面是定义并使用名为 \cmd{direction} 的新数据类型的示例。它有四个允许取值。
\begin{lstlisting}[style=pxl]
direction : newtype enum of {
   NORTH,
   SOUTH,
   EAST,
   WEST
};
compass : direction; // Declare compass as type "direction"
compass = EAST;       // Assign a legal value to "compass"
\end{lstlisting}

\subsubsection{enum of...}
\paragraph*{enum of...}

定义一组与其他类型不同的名称（枚举字面量）。PXL 编码标准要求枚举字面量为全大写字符。

例如：
\begin{lstlisting}[style=pxl]
t3: newtype enum of {
   X,
   Y,
   Z,               // Trailing comma is not needed; silently discarded.
};
\end{lstlisting}

另一示例：
\begin{lstlisting}[style=pxl]
direction : newtype enum of {
   NORTH, SOUTH, EAST, WEST
};

compass : direction;                 // Declare compass as type "direction"
compass = EAST;                      // Assign a legal value to "compass"
\end{lstlisting}

枚举可与其他枚举类型共享名称而无错误：
\begin{lstlisting}[style=pxl]
t1_e : newtype enum of {
   NONE, TOP, BOTTOM, ALL
};
\end{lstlisting}

% ---- 第8章后半 ----
% 第8章后半（合并进 chapters/08_pxl.tex）
% IC Validator User Guide X-2025.06 第 8 章 PXL 后半（p211--243）
% 章节正文片段；不含 \documentclass / preamble
% 续：User-Defined Structures and Types（enum of... 跨类型说明起）

不允许跨类型（cross-type）操作。此外，在同一类型定义内，同一名称不能使用超过一次。

下例展示跨类型操作。该操作不合法：
\begin{lstlisting}[style=pxl]
v1 : t1_e;           // NONE, TOP, BOTTOM, ALL
v1 = NONE;           // Legal assignment
v1 = TOP;            // Legal assignment
v1 = RIGHT;          // Illegal assignment

v2 : t2_e;           // NONE, LEFT, RIGHT, ALL
v2 = NONE;           // Legal assignment

if (v2 == v1) // Illegal
\end{lstlisting}

\subsubsection{\cmd{list of...}}

列表（list）是对象的数组。
\begin{itemize}
  \item 所有对象属于同一类型。该类型允许为空，也可包含重复对象。
  \item 可用 \cmd{foreach()} 循环结构进行处理。
  \item 列表可用方括号（\cmd{[]}）按索引号访问。
\end{itemize}

\begin{seeAlsoBox}
关于列表的列表，见第~\pageref{sec:pxl-nested-containers}~页「嵌套容器」。
\end{seeAlsoBox}

该类型支持表~\ref{tab:pxl-list-ops} 所示操作。

\begin{longtable}{@{}>{\raggedright\arraybackslash\ttfamily}p{0.28\textwidth}
  >{\raggedright\arraybackslash}p{0.12\textwidth}
  >{\raggedright\arraybackslash}p{0.12\textwidth}
  >{\raggedright\arraybackslash}p{0.10\textwidth}
  p{0.28\textwidth}@{}}
\caption{列表操作}\label{tab:pxl-list-ops}\\
\toprule
操作 & e2 类型 & 返回 & 修改 e1 & 说明 \\
\midrule
\endfirsthead
\caption[]{列表操作（续）}\\
\toprule
操作 & e2 类型 & 返回 & 修改 e1 & 说明 \\
\midrule
\endhead
\bottomrule
\endfoot
e1.empty() & -- & boolean & 否 & 若列表不含任何元素则返回 true，否则返回 false。 \\
e1.size() & -- & integer & 否 & 返回列表中的元素个数。 \\
e1.contains(e2) & item & boolean & 否 & 若 e2 是 e1 列表的元素则返回 true。 \\
e1.front() & -- & item & 否 & 返回列表第一个元素的副本。 \\
e1.back() & -- & item & 否 & 返回列表最后一个元素的副本。 \\
e1.chop\_back() & -- & list & 否 & 返回去掉最后一个元素后的 e1 列表副本。 \\
e1.chop\_front() & -- & list & 否 & 返回去掉第一个元素后的 e1 列表副本。 \\
e1.append(e2) & item & list & 否 & 返回将 e2 追加到列表末尾后的 e1 列表副本。 \\
e1.prepend(e2) & item & list & 否 & 返回将 e2 插入到列表前端后的 e1 列表副本。 \\
e1.push\_front(e2) & item & void & 是 & 将 e2 表达式的副本添加到列表前端。 \\
e1.push\_back(e2) & item & void & 是 & 将 e2 表达式的副本添加到列表末尾。 \\
e1.pop\_front() & -- & void & 是 & 移除列表的第一个元素。 \\
e1.pop\_back() & -- & void & 是 & 移除列表的最后一个元素。 \\
e1.merge(e2) & list & list & 否 & 返回将 e2 列表拼接在末尾后的 e1 列表副本。 \\
e1.remove(e2) & item & void & 是 & 移除 e1 列表中所有等于 e2 表达式的值。 \\
\end{longtable}

下列为合法代码示例：
\begin{lstlisting}[style=pxl]
x: layer = ... ;

MyPartition: list of layer = PartitionLayer(L1);

foreach (part in MyPartition) { ... }

if (MyPartition.contains(x)) { ... }

i : integer = MyPartition.size();
\end{lstlisting}

下例展示如何拼接列表：
\begin{lstlisting}[style=pxl]
list1 : list of integer = { 1, 2, 3 };
list2 : list of integer = { 10, 20, 30 };
list3 : list of integer = list1.merge(list2);
\end{lstlisting}

此时 \cmd{list3} 的值为 \cmd{\{ 1, 2, 3, 10, 20, 30 \}}。

下面是初始化和使用字符串列表的示例：
\begin{lstlisting}[style=pxl]
// declare myList to be a list of strings, uninitialized
myList : list of string;

// declare myList to be a list of strings, initialized
myList : list of string = { "this", "is" , "a", "list" };
// myList[0] = "this"
// myList[3] = "list"
\end{lstlisting}

支持内置列表运算符函数。例如：
\begin{lstlisting}[style=pxl]
myList : list of string = { "this", "is" , "a", "list" };

s = myList.size();                 // returns the size of a list: s is "4"
a = myList.front();                // a contains a copy of L1[0]: a is "this"
myList.pop_front();                // removes the first element of the list,
                                   // the list is now { "is", "a", "list" }
myList.pop_back();                 // removes the last element of the list,
                                   // the list is now { "is", "a" }
myList.push_front("xyz");          // adds a new element containing "xyz"
                                   // to the front of the list.
                                   // the list is now { "xyz", "is", "a" }
myList.push_back("bar");           // adds a new element containing "bar"
                                   // to the back of the list.
                                   // the list is now { "xyz", "is", "a", "bar" }
\end{lstlisting}

\subsubsection{\cmd{hash of...}}

定义从一种数据类型到另一种数据类型的映射。源类型称为键类型（key type）；目标类型称为值类型（value type）。也就是说，hash 定义了键/值对数组。

该类型要求键类型为 \cmd{string}、\cmd{integer} 或 \cmd{enum}，或由这三种类型之一派生的类型。不能是 \cmd{double}。值类型可以是任何先前已定义的类型。

下例创建将整数转换为字符串的 hash。此时键类型为 \cmd{integer}，值类型为 \cmd{string}。该类型使用 \cmd{=>} 运算符初始化。
\begin{lstlisting}[style=pxl]
t1_h : newtype hash of integer to string;
t2_h : newtype hash of string to integer;

h1 : t1_h = {
   1 => "a",
   2 => "b",
   26 => "z"
};

h2 : t2_h = {
   "xyzzy" => 1,
   "plugh" => 10,
   "help" => 100
};
\end{lstlisting}

下面是使用派生类型的示例：
\begin{lstlisting}[style=pxl]
t1 : newtype integer;
h1_h : newtype hash of t1 to string;
x : h1_h = { 1 => "a" } ;
\end{lstlisting}

\cmd{hash of...} 类型可使用下标按元素赋值：
\begin{lstlisting}[style=pxl]
x = h1[203];
h1[30] = "ad";
\end{lstlisting}

\cmd{hash of...} 类型支持表~\ref{tab:pxl-hash-ops} 所示成员操作。

\begin{longtable}{@{}>{\raggedright\arraybackslash\ttfamily}p{0.28\textwidth}
  >{\raggedright\arraybackslash}p{0.14\textwidth}
  >{\raggedright\arraybackslash}p{0.12\textwidth}
  >{\raggedright\arraybackslash}p{0.10\textwidth}
  p{0.26\textwidth}@{}}
\caption{Hash 操作}\label{tab:pxl-hash-ops}\\
\toprule
操作 & e2 类型 & 返回 & 修改 e1 & 说明 \\
\midrule
\endfirsthead
\caption[]{Hash 操作（续）}\\
\toprule
操作 & e2 类型 & 返回 & 修改 e1 & 说明 \\
\midrule
\endhead
\bottomrule
\endfoot
e1.contains\_key(e2) & hash 的键类型 & boolean & 否 & 若 e2 是 e1 哈希表的键则返回 true。效果上与 \cmd{e1.keys().contains(e2)} 相同，但运行时性能可能更好。 \\
e1.contains\_value(e2) & hash 的值类型 & boolean & 否 & 若 e2 是 e1 哈希表中某项的值则返回 true。效果上与 \cmd{e1.values().contains(e2)} 相同，但运行时性能可能更好。 \\
e1.empty() & -- & boolean & 否 & 若 hash 不含任何元素则返回 true，否则返回 false。 \\
e1.keys() & -- & list of 键类型 & 否 & 返回所有键值的列表；该列表不含重复值。 \\
e1.size() & -- & integer & 否 & 返回 hash 中的元素个数。 \\
e1.values() & -- & list of 值类型 & 否 & 返回所有值的列表；根据 hash 内容，该列表可含重复值。 \\
e1.merge(e2) & 与 e1 匹配的 hash 类型 & hash & 否 & 返回将 e2 的键与值拼接到末尾后的 e1 hash 副本；若某键在 e1 与 e2 中均存在，则取其在 e2 中的值。 \\
e1.remove(e2) & item & void & 是 & 从 e1 hash 中移除 e2 键。若 e2 键不存在则无效果。 \\
\end{longtable}

再举一例：
\begin{lstlisting}[style=pxl]
alpha_h : newtype hash of string to integer;
alpha : alpha_h = {       // declare and initialize key => value pairs
        "a" => 1,
        "b" => 2,
        "c" => 3
        };
y = alpha["b"];          // The value of y is "2"
keyList = alpha.keys(); // A list of the keys in alpha
\end{lstlisting}

\subsubsection{\cmd{struct of...}}

定义由用户提供名称标识的独立类型。必须至少有一个成员，且成员名必须唯一。各成员可以属于不同的类型。

下面是一个用几何坐标定义点的结构示例：
\begin{lstlisting}[style=pxl]
point: newtype struct of { x: integer; y: integer; };
\end{lstlisting}

\cmd{struct of...} 类型可包含初始值，用作结构创建时的默认值。
\begin{lstlisting}[style=pxl]
t3_s : newtype struct of {
   x: integer;
   y: integer = 76;   // initializer
   z: string = "hi"; // initializer
};
t4_s : newtype struct of {
   a: integer = 106; // initializer
   b: string = "boo"; // initializer
};
\end{lstlisting}

\cmd{struct of...} 类型可从 C 风格初始化器赋值，也可从名--值对列表赋值。下例展示使用上例所定义类型的赋值。
\begin{lstlisting}[style=pxl]
s3 : t3_s = { 10, z = "bye" };            // note that 'y' defaults
                                          // to the value 76

// the defaults cause s4 to have a=106, b="boo"
s4 : t4_s = { };
\end{lstlisting}

再举一例：
\begin{lstlisting}[style=pxl]
shoebox : newtype struct of {
      style : string = "Tennis Shoe";
      shoe_size : double = 10.5;
   };

box : shoebox;
shoe = box.style;          // shoe is "Tennis shoe", the default
box.style = "Dress Shoe"; // Reassign the member box.style
// create a new shoebox with a different kind of shoe
box2 : shoebox = {"Boat Shoe", 9};
\end{lstlisting}

\subsubsection{\cmd{function}}

允许将函数传递给另一个用户函数。被传递的函数可像全局函数一样使用。例如：
\begin{lstlisting}[style=pxl]
// input file
#include <icv.rh>
xyzzy : function (void) returning void { note("abcd"); }
plugh : function ( f : function(void) returning void ) returning void {
f(); }
plugh(xyzzy); // summary file shows "abcd" in output when plugh runs
\end{lstlisting}

\subsubsection{嵌套类型}
\subsubsection*{Nested Types}

类型可以组合形成嵌套类型，例如：
\begin{itemize}
  \item \cmd{list of struct}
  \item \cmd{list of enum}
\end{itemize}

例如，使用上文 \cmd{struct of...} 中 \cmd{shoebox} 结构定义：
\begin{lstlisting}[style=pxl]
box : list of shoebox = {
         { "Walking", 12 },
         { "High Heel", 7 }
      };

A = box[0].style;                  // A is "Walking"
B = box[1].shoe_size;              // B is 7
\end{lstlisting}

\subsection{类型分析}
\subsubsection*{Type Analysis}

本节说明 PXL 在特殊情况下如何处理数据类型，例如当参数传入一种数据类型但期望另一种数据类型时。

当 PXL 遇到一种数据类型而期望另一种时，只要合适就会执行自动类型转换。用于自动类型转换的规则在下列各小节中讨论。

\subsubsection{标准算术转换}
\subsubsection*{Standard Arithmetic Conversions}

当表达式所要求的算术类型与实际找到的算术类型不匹配时，按表~\ref{tab:pxl-arith-conv} 进行转换。

\begin{longtable}{@{}p{0.26\textwidth}
  >{\centering\arraybackslash}p{0.14\textwidth}
  >{\centering\arraybackslash}p{0.16\textwidth}
  >{\centering\arraybackslash}p{0.14\textwidth}
  >{\centering\arraybackslash}p{0.16\textwidth}@{}}
\caption{算术转换}\label{tab:pxl-arith-conv}\\
\toprule
期望 \ldots & 允许 integer\textsuperscript{a} & 允许 constraint of integer\textsuperscript{b} & 允许 double & 允许 constraint of double \\
\midrule
\endfirsthead
\caption[]{算术转换（续）}\\
\toprule
期望 \ldots & 允许 integer & 允许 constraint of integer & 允许 double & 允许 constraint of double \\
\midrule
\endhead
\bottomrule
\endfoot
integer & 是 & 否 & 否 & 否 \\
constraint of integer & 是 & 是 & 否 & 否 \\
double & 是 & 否 & 是 & 否 \\
constraint of double & 是 & 是 & 是 & 是 \\
\end{longtable}

{\footnotesize
\textsuperscript{a}「是」表示允许该转换。\\
\textsuperscript{b}「否」表示不允许该转换。该操作将导致错误。
}

这些转换仅在表达式类型要求时才应用。特别地，除非上下文要求，否则 \cmd{integer} 类型的表达式不会转换为 \cmd{double}。

将 \cmd{integer} 或 \cmd{double} 转换为 constraint 时，所得 constraint 看起来如同使用 \cmd{==} 一元约束（\cmd{CONSTRAINT\_EQ}）书写一样。
\begin{lstlisting}[style=pxl]
f: function ( c: constraint ) returning void { ... }

f( (1.0, 26.0) );                     // constraint is ok
f( 2.0 );                             // single value is ok also
f( 42 );                              // single integer is accepted
\end{lstlisting}

\subsubsection{比较 double}
\subsubsection*{Comparing Doubles}

由于类型 \cmd{double} 是实数的近似，偶尔会出现累积误差，使表面上等价的运算产生不同结果。PXL 提供特殊函数，用于在考虑此类微小误差的情况下比较 double。这些特殊函数见 IC Validator Reference Manual 中「Utility Functions」一章 Math Functions 节的 Equivalence Comparison Functions 与 Double Comparison Functions。

\begin{noteBox}
只要涉及与 double 相关类型之一的直接比较（\cmd{==} 或 \cmd{!=}），就会报告警告。
\end{noteBox}

\subsubsection{类型提升}
\subsubsection*{Type Promotion}

当期望列表，而找到的表达式类型与该列表所含元素类型相同时，该表达式会从单个项提升为仅含该项的列表。

\subsubsection{参数传递}
\subsubsection*{Parameter Passing}

带有 \cmd{output} 或 \cmd{input-output} 限定的函数参数，其实参必须匹配。前述规则不适用。例如，当函数期望标记为 \cmd{out} 的 constraint 参数时，提供 \cmd{integer} 类型变量名不合法，因为所需转换不合法。

\subsubsection{字符串拼接}
\subsubsection*{String Concatenation}

表~\ref{tab:pxl-str-concat} 所示类型的任何表达式，均可使用标准加法运算符（\cmd{+}）追加到字符串表达式。

\begin{longtable}{@{}>{\raggedright\arraybackslash\ttfamily}p{0.32\textwidth} p{0.60\textwidth}@{}}
\caption{字符串拼接}\label{tab:pxl-str-concat}\\
\toprule
类型 & 追加到结果字符串的内容 \\
\midrule
\endfirsthead
\caption[]{字符串拼接（续）}\\
\toprule
类型 & 追加到结果字符串的内容 \\
\midrule
\endhead
\bottomrule
\endfoot
integer & 整数表达式的值 \\
double & double 表达式的值 \\
handle & handle 的文本表示；每个不同对象唯一 \\
string & 字符串内容 \\
boolean & \cmd{true} 或 \cmd{false} \\
constraint of integer & 等价的 constraint 字面量表达式 \\
constraint of double & 等价的 constraint 字面量表达式 \\
user-defined enum & 该表达式所表示的 enum 字面量名称 \\
\end{longtable}

使用字符串时请记住：
\begin{itemize}
  \item 当加法运算符（\cmd{+}）左侧为字符串时，结果为字符串。
  \item 对字符串值表达式，加法运算符（\cmd{+}）不对称：\cmd{a + b != b + a}。
  \item 后续追加到字符串的各项之间不添加空格，因此可能看到意外结果，如下例：
\begin{lstlisting}[style=pxl]
s : string;
s = "hello " + 10 + " " + (1 == 0);
// s now contains the string "hello 10 false"

c : constraint = ( 10, 40 ];
s = "" + c;
// s now contains the string "(10.,40.]", because the
// type of "c" is constraint of double

mye : newtype enum of { A, B, C, D, X };
t : mye = X;
s = "" + t + A + B;
// s now contains the string "XAB"

s = "" + 11 + 3.14 + (1 == 1);
// s now contains the string "113.14true"

s = "" + (11 + 3.14) + (1 == 1);
// s now contains the string "14.14true"
\end{lstlisting}
\end{itemize}

\subsubsection{默认值注入}
\subsubsection*{Default Injection}

处理函数调用时，未给出值的参数若有默认值，则采用函数原型中的默认值。既无默认值也未指定值的参数会导致错误。

类似地，处理字面量结构时，未给出值的成员若有默认值，则采用 \cmd{newtype} 结构定义中的默认值。既无默认值也未指定值的成员会导致错误。

\subsection{嵌套容器}
\subsubsection*{Nested Containers}
\label{sec:pxl-nested-containers}

列表、hash 与结构也称为容器（container）。容器可以多种不同组合嵌套。嵌套容器的一个例子是结构的列表。

下面是 hash 的 hash 示例。runset 代码为（hash of a hash）：
\begin{lstlisting}[style=pxl]
hash1_h : newtype hash of string to integer;
hash2_h : newtype hash of string to hash1_h;

multi_hash : hash2_h = {};

multi_hash["A"] = {};                 // Initialize hash of "A" to empty
multi_hash["A"]["ONE"] = 1;
multi_hash["A"]["TWO"] = 2;
multi_hash["B"] = {
   "THREE" => 3,
   "FOUR" => 4
};

foreach(key1 in multi_hash.keys() ){
   foreach(key2 in multi_hash[key1].keys()){
      note( "multi_hash[" + key1 + "][" + key2 + "] = "
            + multi_hash[key1][key2] + "\n");
   }
}
\end{lstlisting}

输出为：
\begin{lstlisting}
runset.rs:123 multi_hash[A][ONE] = 1
runset.rs:123 multi_hash[A][TWO] = 2
runset.rs:123 multi_hash[B][FOUR] = 4
runset.rs:123 multi_hash[B][THREE] = 3
\end{lstlisting}

下面是列表的列表示例。runset 代码为（list of a list）：
\begin{lstlisting}[style=pxl]
list1_l : newtype list of integer;
list2_l : newtype list of list1_l;

list2 : list2_l = {
   { 0, 1, 2, 3 },
   { 10, 11, 12, 13 }
   };

list2.push_back({ 20, 21, 22, 23 });
list2[1].push_back( 14 );

for( j=0 to (list2.size() - 1) ){
   for( i=0 to (list2[j].size() - 1) ){
      note( "list2[" + j + "][" + i + "] = " + list2[j][i] + "\n");
   }
}
\end{lstlisting}

输出为：
\begin{lstlisting}
runset.rs:143 list2[0][0] = 0
runset.rs:143 list2[0][1] = 1
runset.rs:143 list2[0][2] = 2
runset.rs:143 list2[0][3] = 3
runset.rs:143 list2[1][0] = 10
runset.rs:143 list2[1][1] = 11
runset.rs:143 list2[1][2] = 12
runset.rs:143 list2[1][3] = 13
runset.rs:143 list2[1][4] = 14
runset.rs:143 list2[2][0] = 20
runset.rs:143 list2[2][1] = 21
runset.rs:143 list2[2][2] = 22
runset.rs:143 list2[2][3] = 23
\end{lstlisting}

再举一个列表的列表示例。runset 代码为：
\begin{lstlisting}[style=pxl]
list3 : list of list of integer = {};

list3.push_back( { 0, 1, 2, 3 } );
list3.push_back( { 4, 5, 6 } );
list3.push_back( { 7, 8 } );
list3.push_back( { 9 } );

for( j=0 to (list3.size() - 1) ){
   for( i=0 to (list3[j].size() - 1) ){
     note( "list3[" + j + "][" + i + "] = " + list3[j][i] + "\n");
   }
}
\end{lstlisting}

输出为：
\begin{lstlisting}
runset.rs:159 list3[0][0] = 0
runset.rs:159 list3[0][1] = 1
runset.rs:159 list3[0][2] = 2
runset.rs:159 list3[0][3] = 3
runset.rs:159 list3[1][0] = 4
runset.rs:159 list3[1][1] = 5
runset.rs:159 list3[1][2] = 6
runset.rs:159 list3[2][0] = 7
runset.rs:159 list3[2][1] = 8
runset.rs:159 list3[3][0] = 9
\end{lstlisting}

% ============================================================
\section{流程控制}
\subsection*{Flow Control}

流程可用下列方式控制：
\begin{itemize}
  \item 条件语句（Conditional Statements）
  \item 循环（Loops）
  \item 执行终止（Termination of Execution）
\end{itemize}

\subsection{条件语句}
\subsubsection*{Conditional Statements}

PXL 中的条件语句是一类构造，允许你根据指定条件求值为 true 还是 false 来执行不同动作。

\cmd{if} 语句用于仅在预定义条件满足时执行某动作。

\cmd{if} 语句至少与另外两个组成部分一起使用：
\begin{itemize}
  \item 一个可求值为布尔值 true 或 false 的表达式。
  \item 一个在求值表达式返回 true 时执行的组成部分。
\end{itemize}

可选地，在必需组成部分之外，还可包含下列组成部分以构建更复杂的条件语句：
\begin{itemize}
  \item \cmd{else} 构造，后跟在第一个表达式求值为 false 时执行的组成部分。
  \item \cmd{elseif} 构造，表示在父 \cmd{if} 条件语句内嵌套了循环式 \cmd{if} 条件或一个或多个 \cmd{if} 条件。PXL 也支持 \cmd{elif} 构造。
\end{itemize}

\cmd{if} 语句使用下列语法：
\begin{lstlisting}[style=pxl]
if (e1) s1;
if (e1) s1 elif (e2) s2;
if (e1) s1 elseif (e2) s2;
if (e1) s1 else s2;
\end{lstlisting}

\cmd{if} 语句按如下方式求值：
\begin{itemize}
  \item 始终先求值第一个布尔表达式（\cmd{e1}）。
  \item 若 \cmd{e1} 求值为 true，则处理第一个组成部分（\cmd{s1}）。
  \item 若遇到 \cmd{else} 语句，则执行其后的语句（\cmd{s2}）。
  \item 若遇到 \cmd{elif} 或 \cmd{elseif} 语句，则执行其后的语句。
\end{itemize}

\subsection{循环}
\subsubsection*{Loops}

与其他编程语言一样，可在 PXL 中使用循环将某动作重复若干次。

PXL 提供下列方法创建循环构造：
\begin{itemize}
  \item \cmd{for} 循环
  \item \cmd{foreach} 循环
  \item \cmd{while} 循环
\end{itemize}

\subsubsection{\cmd{for} 循环}
\subsubsection*{for Loop}

\cmd{for} 循环用于将特定动作重复指定次数。它至少包含一个整数（\cmd{e1}）以及必须重复的动作（\cmd{s1}）。也可包含附加整数（\cmd{e2}、\cmd{e3}）。

语法为：
\begin{lstlisting}[style=pxl]
for (var = e1 to e2) s1;         //s1 is a statement; it can be compound
for (var = e1 to e2 step e3) s1;
\end{lstlisting}

\begin{itemize}
  \item 要求 \cmd{e1}、\cmd{e2} 与 \cmd{e3} 为整数值表达式。
  \item 将 \cmd{var} 定义为仅在语句 \cmd{s1} 内可见的整数变量。
  \item 在循环开始时对 \cmd{e1}、\cmd{e2} 与 \cmd{e3} 各求值一次；随后对这些表达式取值的赋值不影响循环。
\end{itemize}

例如：
\begin{lstlisting}[style=pxl]
for(i=0 to 10 step 2){              // Increment by 2: 0, 2, 4, 6, . . .
   if( i < 6) {
      a[i] = b[i] + c[i];
   } elif( ( i >=6 ) && ( i < 10)){
      a[i] = b[i] + c[i] + 1;
   } else {
      a[i] = b[i] - c[i];
   }
};
\end{lstlisting}

\subsubsection{\cmd{foreach} 循环}
\subsubsection*{foreach Loop}

\cmd{foreach} 循环类似于 \cmd{for} 循环，但用于针对容器（如列表或 enum）的元素执行一组动作。它不使用变量计数器来决定动作的迭代次数。

它包含类型为列表的 \cmd{e1}，以及对 \cmd{e1} 中每个元素求值的组成部分 \cmd{s1}。

语法为：
\begin{lstlisting}[style=pxl]
foreach (var in e1) s1;         //e1 is of the type "list of"
\end{lstlisting}

\begin{itemize}
  \item 对 \cmd{e1} 的每个元素求值一次 \cmd{s1}，同时将 \cmd{var} 从列表复制过来。
  \item 若输入列表 \cmd{e1} 为空，则跳过 \cmd{s1}。
  \item \cmd{var} 是类型正确的常量变量，仅在语句 \cmd{s1} 内可见。
\end{itemize}

\subsubsection{\cmd{while} 循环}
\subsubsection*{while Loop}

\cmd{while} 循环用于重复特定动作直到满足某条件。它包含布尔表达式（\cmd{e1}）以及必须重复的动作（\cmd{s1}）。\cmd{while} 循环始终至少对布尔表达式 \cmd{e1} 求值一次。

语法为：
\begin{lstlisting}[style=pxl]
while (e1) s1;       // e1 is evaluated as a Boolean;
                     // loop continues until e1 is false
\end{lstlisting}

\cmd{while} 语句按如下方式求值：
\begin{itemize}
  \item 每次 \cmd{e1} 求值为 true 时，执行 \cmd{s1}。
  \item 当 \cmd{e1} 求值为 false 时，该语句结束。
\end{itemize}

\subsection{执行终止}
\subsubsection*{Termination of Execution}

\cmd{error("message")} 是诊断函数（Diagnostics Function）。它写出所提供的消息，然后终止执行。例如：
\begin{lstlisting}[style=pxl]
#ifdef ABC
#ifdef DEF
error("ABC and DEF are mutually exclusive. Please do not set both.");
#endif
#endif
\end{lstlisting}

% ============================================================
\section{函数}
\subsection*{Functions}

PXL 中的函数按用途分类。函数类型见表~\ref{tab:pxl-func-types}。

\begin{longtable}{@{}>{\raggedright\arraybackslash}p{0.22\textwidth} p{0.70\textwidth}@{}}
\caption{函数类型}\label{tab:pxl-func-types}\\
\toprule
函数类型 & 说明 \\
\midrule
\endfirsthead
\caption[]{函数类型（续）}\\
\toprule
函数类型 & 说明 \\
\midrule
\endhead
\bottomrule
\endfoot
Runset function & IC Validator 发布的标准函数。 \\
Remote function & 由用户编写、可被 runset 函数调用的函数。IC Validator 工具也提供默认 remote 函数。 \\
User function & 由用户编写的函数。 \\
Utility function & 由 IC Validator 工具提供、用于在 remote 函数中访问原始数据的函数。 \\
\end{longtable}

\subsection{函数定义}
\subsubsection*{Function Definitions}

函数定义为一组与名称关联的操作。函数定义标识所有参数以及返回值的名称、类型与默认值。

你可以创建自己的函数，并与所提供的函数一起在 runset 中使用。

如下例所示，所有函数定义必须位于顶层。
\begin{lstlisting}[style=pxl]
// a simple function definition

myinversion: function ( arg: double ) returning rv : double
{
   rv = 1.0 / arg;
}

x = myinversion(2.0);         // simple function call
\end{lstlisting}

函数定义为所有处理提供单独作用域。作用域由书写位置的上下文决定（静态作用域，static scoping）。

\begin{seeAlsoBox}
关于作用域，见「变量作用域」（原书第~197~页）。
\end{seeAlsoBox}

在函数定义内部适用下列作用域规则：
\begin{enumerate}
  \item 若在当前作用域找不到声明，则向上搜索更高作用域，直到搜遍整个函数定义，在匹配处停止。
  \item 然后，如有必要，使用在函数定义之前声明且同名的全局变量。
  \item 最后，若该变量正被赋值为函数调用的结果，或处于函数调用的输出参数位置，则在当前语句的作用域中隐式声明它。但是，若该变量处于函数调用的输入或输入--输出参数位置，即使它正被赋值为函数调用的结果，也不会在当前语句的作用域中隐式声明。
\end{enumerate}

在函数定义外部适用下列作用域规则：
\begin{enumerate}
  \item 若在当前作用域找不到声明，则向上搜索更高作用域，直到搜遍整个 runset，在匹配处停止。
  \item 然后，如有必要，使用先前声明且同名的全局变量。
  \item 最后，若该变量正被赋值为函数调用的结果，或处于函数调用的输出参数位置，则在当前语句的作用域中隐式声明它。但是，若该变量处于函数调用的输入或输入--输出参数位置，即使它正被赋值为函数调用的结果，也不会在当前语句的作用域中隐式声明。
\end{enumerate}

\subsection{函数原型}
\subsubsection*{Function Prototypes}

函数原型定义函数的参数以及可选参数的默认值。函数原型标识按引用调用（call-by-reference）的参数。函数原型不定义函数的操作。

下面是函数原型示例：
\begin{lstlisting}[style=pxl]
my_xor : function (
   a : polygon_layer,
   b : polygon_layer
) returning xorOut: polygon_layer;
\end{lstlisting}

下面是函数声明示例：
\begin{lstlisting}[style=pxl]
my_xor : function (
   a : polygon_layer,
   b : polygon_layer
) returning xorOut: polygon_layer {
      tmp1 = not(a,b);
      tmp2 = not(b,a);
      xorOut = or(tmp1, tmp2 );
   };
\end{lstlisting}

下例展示调用前一例中定义的 \cmd{my\_xor()} 函数：
\begin{lstlisting}[style=pxl]
a1 : polygon_layer;
b1 : polygon_layer;
xorA1B1 = my_xor(a1, b1);
\end{lstlisting}

\subsection{函数调用}
\subsubsection*{Function Calls}

函数调用定义函数如何运作。函数调用由函数名后跟通过实参说明如何及在何处执行该函数的指令组成。

向函数调用的实参传递值可用两种不同风格：
\begin{itemize}
  \item 仅按实参值（按值传递，pass-by-value）。值的顺序必须与函数原型中参数顺序匹配。使用此风格时，若你关心某个可选参数，则必须为其前面的所有可选参数指定值，因为值按顺序匹配到参数。
  \item 按 \cmd{argument\_name = argument\_value} 格式用名称（按名传递，pass-by-name）。参数顺序无关。使用此风格时，无需为其前面的所有可选参数指定值，因为值按名称匹配到参数。
\end{itemize}

两种风格可在同一次函数调用中使用。但是，在函数调用中，所有按值传递风格指定的实参值必须出现在所有按名传递风格指定的实参值之前。一旦在函数调用中使用了按名传递风格，其后的所有实参调用必须使用相同风格。

例如，\cmd{interacting()} 函数的语法为：
\begin{lstlisting}[style=pxl]
interacting : binary published function(
   layer1 : polygon_layer,
   layer2 : polygon_layer,
   count : count_constraint_t = >0,
   include_touch : more_include_touch_e = EDGE,
   processing_mode : processing_mode_e = HIERARCHICAL
) returning result : polygon_layer;
\end{lstlisting}

下面是以按名传递风格调用 \cmd{interacting()} 的示例：
\begin{lstlisting}[style=pxl]
y = interacting( layer1 = a, layer2 = b, include_touch = ALL );
   // The count argument is skipped, and the value is the default.
   // The value of the processing_mode argument is the default.
\end{lstlisting}

下面是以按值传递风格调用 \cmd{interacting()} 的示例：
\begin{lstlisting}[style=pxl]
y = interacting( a, b, 2 );
   // The value of count is 2.
   // The values of the include_touch and processing_mode arguments
   // are the defaults.
\end{lstlisting}

下面是两种混合风格调用 \cmd{interacting()} 的示例：
\begin{lstlisting}[style=pxl]
y = interacting( a, b, include_touch = ALL );
   // The count argument is skipped, and the value is the default.
   // The value of the processing_mode argument is the default.

y = interacting( a, b, count = 2);
   // The values of the include_touch and processing_mode arguments
   // are the defaults.
\end{lstlisting}

函数调用有两种风格：
\begin{itemize}
  \item 标准函数调用（Standard Function Call）
\begin{lstlisting}[style=pxl]
y = and(a, b);
y = and(a, b, CELL_LEVEL);
y = or(not(a,b),not(b,a));
\end{lstlisting}
  \item 运算符风格函数调用（Operator-Style Function Call）
\begin{lstlisting}[style=pxl]
y = a and b;
y = a and[CELL_LEVEL] b;
y = (a not b) or (b not a);
\end{lstlisting}
\end{itemize}

\subsubsection{标准函数调用}
\subsubsection*{Standard Function Call}

对任何函数始终允许标准函数调用。

标准函数调用由函数名后跟括号括起、逗号分隔的至少一个表达式列表，再跟至少一个名--值对列表组成。

下例中函数有三个参数：
\begin{lstlisting}[style=pxl]
g: function ( x : integer, y : integer, z : integer)
   returning void;
\end{lstlisting}

下列四个函数调用中，有一个不合法：
\begin{lstlisting}[style=pxl]
g(x=3, y=10, z=10);          // allowed

g(z=3, y=10, x=4);           // allowed

g(10, 20, 30);               // allowed; x=10, y=20, z=30

g(z=3, 10, 20);              // ILLEGAL - expressions must not
                             // follow name-value pairs
\end{lstlisting}

若任一参数存在默认值，则不必指定这些参数。下例中函数有三个参数，其中两个有默认值。
\begin{lstlisting}[style=pxl]
f: function
    ( x : integer,
      y : integer = 27,
      z : integer = 42
    ) returning void { ... }
...
\end{lstlisting}

下列五个函数调用中，有两个不合法：
\begin{lstlisting}[style=pxl]
f(10, y=2);          // allowed, z is 42

f(10);               // allowed, y is 27, z is 42

f(z=3, x=4);         // allowed, y is 27

f();                 // ILLEGAL - x requires a value

f(z=20);             // ILLEGAL - x requires a value
\end{lstlisting}

\subsubsection{运算符风格函数调用}
\subsubsection*{Operator-Style Function Call}

运算符风格函数调用要求二元函数在函数名左右两侧各有实参。其他实参须以名--值对形式写在紧随函数名之后的方括号中。

仅对以 \cmd{binary} 限定符声明或定义的函数允许运算符风格函数调用。使用运算符风格函数调用有助于防止在其他上下文（如变量名与成员名）中使用该函数名。

下列语法定义如何使用 \cmd{binary} 限定符调用函数：
\begin{lstlisting}[style=pxl]
length: binary function (l1: layer, c1: constraint of double) returning
result: layer {...}
...
LONGMET = MET1 length > 5;
\end{lstlisting}

通过仔细的接口定义，\cmd{binary} 可写出易读的程序。例如：
\begin{lstlisting}[style=pxl]
// A simple declaration that takes two arguments.
// The LHS is the layer to be selected from; the RHS is the
// constraint defining the criteria area.
...
// ... as is a variable of appropriate type.
c : constraint = (27.0, 42.5);
x : lyr1 area c;
\end{lstlisting}

\subsubsection{实参绑定}
\subsubsection*{Argument Bindings}

实参绑定决定调用点的哪些实参绑定到函数原型中的哪些参数。
\begin{itemize}
  \item 根据函数名以及上下文所要求的类型选择候选原型。若无法确定类型，则仅使用函数名。进行运算符风格函数调用时，外部实参绑定到原型中的参数。对二元函数，左侧实参绑定到第一个参数，右侧实参绑定到第二个参数。
  \item 每个名--值实参绑定到剩余的同名参数。
  \item 任何没有名称的实参从左到右绑定到剩余的无默认值参数。
  \item 若所有剩余未绑定参数都有默认值，则使用该函数。
  \item 若实参绑定未能恰好匹配一个函数原型，则为错误。
\end{itemize}

例如，给定如下定义的函数 \cmd{MyFunc}：
\begin{lstlisting}[style=pxl]
MyFunc: binary function (
   fromLayer: lhs layer,
   touching: rhs layer,
   corners: boolean = true,
   whole_polygon: boolean
) returning e: edge_layer;
\end{lstlisting}

下列函数调用全部等价：
\begin{lstlisting}[style=pxl]
T1 = MyFunc (metalLayer, selectLayer, true);
T2 = MyFunc (metalLayer, selectLayer, whole_polygon=true);
T3 = metalLayer MyFunc[whole_polygon=true] selectLayer;
\end{lstlisting}

下列代码因 \cmd{whole\_polygon} 实参没有默认值而产生编译时错误：
\begin{lstlisting}[style=pxl]
/* illegal function call */
Tbroken = metalLayer MyFunc selectLayer;
\end{lstlisting}

\subsubsection{默认值}
\subsubsection*{Defaults}

默认值在函数原型中指定为选项变量的右侧赋值。例如，\cmd{and()} 函数的语法为：
\begin{lstlisting}[style=pxl]
and : function(
   layer1 : polygon_layer,
   layer2 : polygon_layer,
   processing_mode : processing_mode_e = HIERARCHICAL,
   remove_hierarchical_overlap : boolean = true
) returning and_result : polygon_layer;
\end{lstlisting}

下列对 \cmd{and()} 的调用对 \cmd{processing\_mode} 与 \cmd{remove\_hierarchical\_overlap} 实参使用默认值：
\begin{lstlisting}[style=pxl]
c = and( a, b );
\end{lstlisting}

下列对 \cmd{and()} 的调用对 \cmd{processing\_mode} 使用可选设置，对 \cmd{remove\_hierarchical\_overlap} 使用默认值：
\begin{lstlisting}[style=pxl]
c = and( a, b, CELL_LEVEL );
\end{lstlisting}

\paragraph{重置默认值}
\paragraph*{Resetting Defaults}

你可以覆盖任何函数实参的默认值。任何实参都可更改其默认值；可以是数值、布尔值或枚举量。

本例中设置了 \cmd{abc} 函数实参 \cmd{x} 的默认值：
\begin{lstlisting}[style=pxl]
abc : literal function ( x : double = 123) returning z : double;
A = abc();   // == abc(123)
\end{lstlisting}

\begin{noteBox}
警告：重置默认值时，你是在重新定义该函数，必须列出其全部实参。
\end{noteBox}

\subsection{匿名函数}
\subsubsection*{Anonymous Functions}

若干 IC Validator 命令允许将用户函数作为实参传入。匿名函数为简单行为或组织更复杂行为提供了一种机制。可在期望函数名的地方使用匿名函数。编译器生成内部名称，并在紧临其调用之前的代码位置声明该函数。

传统上，可使用全局变量，并按需为不同调用设置不同值。例如：
\begin{lstlisting}[style=pxl]
x : integer;
y : integer;
calc_nmos_props : function (void) returning void { ... }

x = 0;
y = 5;
nmos(..., calc_nmos_props, ...);

x = 1;
y = 4;
nmos(..., calc_nmos_props, ...);
\end{lstlisting}

然而，匿名函数可使此类代码更紧凑、组织更好。下例中，在每次 NMOS 调用之前声明一个隐藏函数，随后使用它。
\begin{lstlisting}[style=pxl]
calc_nmos_props : function (x : integer, y : integer) returning void
 { ... }

nmos(..., function { calc_nmos_props(0, 5); }, ...);
nmos(..., function { calc_nmos_props(1, 4); }, ...);
\end{lstlisting}

\subsection{函数重载}
\subsubsection*{Function Overloading}

在某些情况下，PXL 允许存在多个同名函数。若 PXL 能根据类型与参数名确定要调用哪个函数，则多个函数可以同名。若 PXL 无法确定要调用的函数，则导致错误。

若一个返回值而另一个返回 \cmd{void}，则可以有两个同名的不同函数。

当函数出现在语句级时，假定为返回 \cmd{void} 的函数；当出现在表达式中时，假定为返回值的函数：
\begin{lstlisting}[style=pxl]
MyFunction(a,b,c);             // use the void-returning function
                               // "MyFunction"
x = MyFunction(a,b,c);         // use the value-returning function
                               // "MyFunction"
\end{lstlisting}

\begin{noteBox}
PXL 仅支持基于不同返回类型的函数重载。
\end{noteBox}

\subsection{函数递归}
\subsubsection*{Function Recursion}

PXL 不支持递归函数调用。例如：
\begin{lstlisting}[style=pxl]
noCanDo : function( ... ) returning integer {
  ...
  x = noCanDo(...);
  ...
}
\end{lstlisting}

% ============================================================
\section{程序结构}
\subsection*{Program Structure}

PXL 程序是声明、函数定义与语句的序列。程序自上而下执行。

PXL 对程序有效性设定下列前提条件：
\begin{itemize}
  \item 所有对象必须在使用前定义。
  \item 所有语句仅因副作用而求值。表达式必须是赋值的一部分。
\begin{lstlisting}[style=pxl]
10 + 5;     // illegal because the result (15) is not captured
\end{lstlisting}
\end{itemize}

\subsection{复合语句}
\subsubsection*{Compound Statement}

复合语句定义可声明局部变量的嵌套作用域。复合语句可用在允许单条语句的任何地方。

\begin{seeAlsoBox}
关于作用域，见「变量作用域」（原书第~197~页）。
\end{seeAlsoBox}

\begin{itemize}
  \item 复合语句允许局部变量隐藏其他作用域级别先前定义的变量。编译器报告警告。隐藏也称为遮蔽（shadowing）。
  \item 复合语句使用花括号（\cmd{\{\}}）标识。
\begin{lstlisting}[style=pxl]
{ ; }      // Smallest legal compound statement
\end{lstlisting}
  \item 复合语句遵循下列语法：
\begin{lstlisting}[style=pxl]
{
    s1;
    s2;
    ...
    sN;
}
\end{lstlisting}
\end{itemize}

\subsection{违例}
\subsubsection*{Violations}

某些 IC Validator 函数可产生违例数据（也称为错误输出），替代或补充层输出。这些违例存储在 SQLite 数据库（错误数据库，PYDB）中，可供 VUE 用于调试，并在运行结束时写入 \cmd{LAYOUT\_ERRORS} 文件。也有一些 PXL 函数以违例数据为输入，例如 \cmd{write\_gds()} 与 \cmd{violation\_empty()} 函数。

若希望直接访问 IC Validator 工具产生的错误数据，见附录~E「PYDB Perl API」，了解如何使用 IC Validator Perl 应用程序编程接口（API）。

\begin{seeAlsoBox}
附录~E，PYDB Perl API。
\end{seeAlsoBox}

\subsubsection{违例注释}
\subsubsection*{Violation Comments}

所有违例输出都与一条注释字符串关联。该注释作为标题出现在 \cmd{LAYOUT\_ERRORS} 文件以及 VUE 中。任意数量的函数可输出到同一违例注释。

当前及嵌套 PXL 作用域的违例注释使用 \cmd{@} 运算符指定。若 runset 中未指定注释，违例获得默认注释 \cmd{"Violation"}。

下例展示如何为层集合中的每一层编写唯一的违例注释。违例命名为 ``Metal1 spacing <= 0.1um''、``Metal2 spacing <= 0.1um''，依此类推。
\begin{lstlisting}[style=pxl]
for (i=0 to metal_layers.size()-1) {
   @ "Metal" + (i+1) + " spacing <= 0.1um";
   external1(metal_layers[i], distance<=0.1);
}
\end{lstlisting}

下例展示违例注释相对于作用域的行为。

\begin{noteBox}
通常，违例注释定义（\cmd{@} 运算符）不应与其所适用的函数分开。下例突出说明违例注释相对于作用域的行为。
\end{noteBox}

\begin{lstlisting}[style=pxl]
@ "First Error Message";
{ @ "Second Error Message";             // A new scope is created here
  local_var = abc(layer2);
  internal1(local_var, distance < 0.2); // Errors from internal1()
                                        // are stored in the database
                                        // under "Second Error Message"
} // The original scope is restored here

external1(layer1, distance < 0.1); // Errors from external1() are stored
                                   // in the database under "First Error
                                   // Message"
\end{lstlisting}

下例展示如何为层集合中的每一层编写唯一的违例注释。违例命名为 ``Metal1 spacing <= 0.1um''、``Metal2 spacing <= 0.1um''，依此类推。
\begin{lstlisting}[style=pxl]
for (i=0 to metal_layers.size()-1) {
   @ "Metal" + (i+1) + " spacing <= 0.1um";
   external1(metal_layers[i], distance<=0.1);
}
\end{lstlisting}

\cmd{text\_net()}、\cmd{texted\_with()} 与 \cmd{not\_texted\_with()} 函数可输出多种不同类型的违例。对这些函数，有办法为个别错误类型覆盖当前 runset 违例注释。本例中，text short 输出到违例注释 ``Rule 4.5.6: No Shorts''，而所有其他类型的 text 错误输出到 ``Rule 1.2.3: No text errors''。
\begin{lstlisting}[style=pxl]
@ "Rule 1.2.3: No text errors";

cdb1 = text_net(cdb1,
                { { M1, M1_text } },
                shorted_violation_comment = "Rule 4.5.6: No shorts");
\end{lstlisting}

IC Validator 工具使用诸如 \cmd{-svc} 等命令行选项处理违例注释。例如：
\begin{lstlisting}[style=pxl]
{@ "A";
   {@ "B";
      internal1(...);
   };
};
\end{lstlisting}

工具在选择时仅考虑较低层的违例注释。在此情况下，设置 \cmd{-svc B} 或 \cmd{-uvc A} 会执行 \cmd{internal1()} 函数；而设置 \cmd{-svc A} 或 \cmd{-uvc B} 则不会。

\subsubsection{产生违例的函数}
\subsubsection*{Violation-Producing Functions}

大多数 DRC 与数据生成函数是重载的；也就是说，它们有两个版本。返回多边形层的版本不产生错误输出。返回 \cmd{void} 的版本产生错误输出但不产生多边形层。下面是 \cmd{external1()} 函数两个版本的示例。
\begin{lstlisting}[style=pxl]
// This external1 outputs to a polygon layer, no errors are produced
t1 = external1(layer1, distance < 0.1);

// This external1 outputs error data to the comment "Ext1 Check"
@ "Ext1 Check";
external1(layer1, distance < 0.1);
\end{lstlisting}

也有一些 PXL 函数不返回 \cmd{void}，但也可产生错误输出。例如，\cmd{text\_net()} 函数返回 connect database，同时也为 text short、open 与 unused text 产生错误输出。

\subsubsection{违例变量}
\subsubsection*{Violation Variables}

内置 PXL 类型 \cmd{violation} 可用于保存错误数据的变量，以便稍后在 runset 中将其用作另一 PXL 函数的输入。例如，违例变量可用作 \cmd{write\_gds()} 等函数的输入以输出错误多边形，或用作 \cmd{violation\_empty()} 以判断违例是否包含任何错误。

如下例所示，某些 PXL 函数直接返回违例：
\begin{lstlisting}[style=pxl]
// These functions return violations for which there is not an associated
// PXL function
v1 = get_layout_grid_violation();
v2 = get_layout_drawn_violation();

// Write errors to GDSII
g = gds_library("err.gds");
write_gds(output_library = g,
          errors = { { v1, { 1 } },
                     { v2, { 2 } } });
\end{lstlisting}

\begin{noteBox}
错误输出可受 \cmd{error\_options()} 函数中 \cmd{error\_limit\_per\_check} 等实参限制。这些选项可影响实际存入错误数据库的数据量。被丢弃的错误对 \cmd{write\_gds()} 及其他函数不可用。
\end{noteBox}

\subsubsection{违例块}
\subsubsection*{Violation Blocks}

违例变量也可使用违例块赋值。违例块使用 \cmd{@=} 运算符指定。左侧为违例变量。\cmd{@=} 运算符后跟花括号括起的代码块。在该块内，所有产生错误输出的函数都与该违例变量关联。例如：
\begin{lstlisting}[style=pxl]
// v1 contains the error output from the external1() check
v1 @= {
   @ "Rule 1";
   l1 = a and b not c;
   external1(l1, distance < 2);
}

// v2 contains the error output from both internal1() checks and
// the external1() check
v2 @= {
   @ "Rule 2";
   internal1(l2, distance < 1);
   internal1(l3, distance < 1);

   @ "Rule 3";
   external1(l2, distance < 1);
}

tdb : connect_database; // Declared outside the violation block scope
                        // so it can be used after the violation block
// v3 contains all text errors found by text_net()
v3 @= {
   @ "Rule 4";
   tdb = text_net(cdb, {{ M1, M1_text }});
}

// These functions can be used to get only certain types of errors
// from a violation. The result is a subset of the input violation.
v4 = get_text_shorted(v3); // v4 contains only text shorts from v3
v5 = get_text_unused(v3); // v5 contains only unused text from v3

// Take some action if v4 is not empty
if (! violation_empty(v4)) {
   note("Found text shorts!");
}
\end{lstlisting}

\subsection{执行模型}
\subsubsection*{Execution Model}

PXL 的执行模型要求资源可用。若某程序中操作的前提条件已满足，则该操作随时可供执行。强烈偏向尽早执行诊断，意味着在程序初始处理期间，所有已就绪的诊断尽可能立即执行。因此，你可能看到 PXL 程序中语句乱序执行。最终结果不受此重排影响；但是，中间结果可能在不同时间出现或同时出现。这种乱序执行也称为并行执行。

同一程序在不同次执行时始终生成相同结果。不过，这些结果可能因资源可用性而以不同顺序出现。可能的运行时错误包括：
\begin{itemize}
  \item 试图使用尚未赋值的任何变量。
  \item 索引超出列表末尾，或以负索引索引。
  \item 以不存在元素的键值从 hash 读取。
\end{itemize}

\subsection{预处理器}
\subsubsection*{Preprocessor}

PXL 使用预处理器使程序的开发与执行更简便、更快速。标准预处理器模块集成在 PXL 编译器中。表~\ref{tab:pxl-preprocessor} 展示预处理器能力。

\begin{noteBox}
预处理器的行长度限制为 64K 字符。超出此限制通常可通过在空白或分隔符（如逗号 \cmd{,}）处插入回车来规避。
\end{noteBox}

\begin{longtable}{@{}>{\raggedright\arraybackslash}p{0.28\textwidth} p{0.64\textwidth}@{}}
\caption{标准预处理器模块}\label{tab:pxl-preprocessor}\\
\toprule
典型预处理能力 & 所用语法 \\
\midrule
\endfirsthead
\caption[]{标准预处理器模块（续）}\\
\toprule
典型预处理能力 & 所用语法 \\
\midrule
\endhead
\bottomrule
\endfoot
条件编译 & \cmd{\#if} \ldots\ \cmd{\#ifdef}、\cmd{\#ifndef}、\cmd{\#elif}、\cmd{\#else} \\
错误生成 & \cmd{\#error} \\
文件包含 & \cmd{\#include} 用于将附加文件内容插入 runset。文件名必须用双引号或尖括号括起。语法控制搜索路径。\newline
用尖括号包含 IC Validator 文件：\newline
\cmd{\#include <icv.rh>}\newline
按下列顺序搜索目录：\newline
1.~命令行上用 \cmd{-I} 选项指定的目录，按其指定顺序\newline
2.~\cmd{ICV\_INSTALL\_DIRECTORY/include}\newline
用双引号包含你自己的文件：\newline
\cmd{\#include "my\_include.rh"}\newline
按下列顺序搜索目录：\newline
1.~含有 \cmd{\#include} 语句的文件所在同一目录\newline
2.~命令行上用 \cmd{-I} 选项指定的目录，按其指定顺序\newline
3.~\cmd{ICV\_INSTALL\_DIRECTORY/include}\newline
指定绝对路径时，仅搜索该路径。\newline
可用 \cmd{\#pragma once} 防止文件被包含超过一次。任何含有 \cmd{\#pragma once} 的文件只会被包含一次，即使它出现在多个 \cmd{\#include} 语句中。\newline
例如，若 runset 目录结构如下：\newline
\cmd{Runset}\newline
\cmd{|-- rules.rs}\newline
\cmd{|-- Include}\newline
\cmd{\phantom{|-- }|-- assign.rs}\newline
\cmd{\phantom{|-- }|-- other\_assign.rs}\newline
要在 \cmd{rules.rs} 中包含 \cmd{other\_assign.rs}，使用\newline
\cmd{\#include "Include/other\_assign.rs"}\newline
要在 \cmd{assign.rs} 中包含 \cmd{other\_assign.rs}，使用\newline
\cmd{\#include "other\_assign.rs"} \\
位置跟踪 & \cmd{\#line} \\
标准宏支持 & \cmd{\#define} \ldots\ \cmd{\#undef}，包括记号构造（宏体中的 \cmd{\#\#} 二元运算符）与字符串化（宏体中的 \cmd{\#} 一元运算符）。也可用 \cmd{-D} 命令行选项定义变量。 \\
\end{longtable}

\begin{seeAlsoBox}
IC Validator 命令行选项（原书第~25~页）。\\
关于 IC Validator 主目录的更多信息，见 SolvNetPlus 上的 IC Validator Installation Notes。
\end{seeAlsoBox}

使用预处理器指令移除大块代码；也就是说，将该区段包在一对 \cmd{\#if 0} 与 \cmd{\#endif} 语句内。例如：
\begin{lstlisting}[style=pxl]
#if 0
// everything through the matching #endif is treated as a comment.
if (something == true)
{
   a line comment.
   do_something(useful = true);
}
#endif
\end{lstlisting}

下面是使用指令强制单次包含的示例。首次遇到此文件时 \cmd{MYFILE\_RH} 未定义。所含的 \cmd{\#define MYFILE\_RH} 将该变量设为已定义，并包含代码体。此后因 \cmd{MYFILE\_RH} 已存在且 \cmd{\#ifndef} 求值为 false，代码体不再被包含。
\begin{lstlisting}[style=pxl]
#ifndef MYFILE_RH
#define MYFILE_RH

// Body of code to include goes here

#endif        // end if the #ifndef MYFILE_RH
\end{lstlisting}

使用 \cmd{\#if} 语句时：
\begin{itemize}
  \item 使用标准比较运算符：\cmd{==}、\cmd{<=}、\cmd{>=}、\cmd{\&\&}、\cmd{||}
  \item 未定义的宏被解释为 0。例如，若 \cmd{SOMEMACRO} 未定义，则 \cmd{\#if SOMEMACRO == 1} 被翻译为 \cmd{\#if 0 == 1}。
  \item 仅比较整数值。不能比较字符串。
\end{itemize}

\subsubsection{预定义标识符}
\subsubsection*{Predefined Identifiers}

预处理期间支持表~\ref{tab:pxl-predef-id} 所示预定义标识符。

\begin{longtable}{@{}>{\raggedright\arraybackslash\ttfamily}p{0.36\textwidth} p{0.56\textwidth}@{}}
\caption{预定义标识符}\label{tab:pxl-predef-id}\\
\toprule
标识符 & 内容 \\
\midrule
\endfirsthead
\caption[]{预定义标识符（续）}\\
\toprule
标识符 & 内容 \\
\midrule
\endhead
\bottomrule
\endfoot
\_\_LINE\_\_ & 当前源文件中的行号，为整数 \\
\_\_FILE\_\_ & 当前源文件名，为字符串常量 \\
\_\_DATE\_\_ & 编译日期，为 11 字符字符串字面量，例如 ``Apr 25 2006'' \\
\_\_TIME\_\_ & 编译时间，为字符串字面量（``hh:mm:ss''） \\
\_\_PXL\_VERSION\_MAJOR\_\_ & 语言主版本号，为整数 \\
\_\_PXL\_VERSION\_MINOR\_\_ & 语言次版本号，为整数 \\
\_\_PXL\_VERSION\_PATCH\_\_ & 语言补丁级别，为整数 \\
\_\_PXL\_DATE\_\_ & 编译日期，为八位整数常量 yyyymmdd，例如 20060425 \\
\_\_PXL\_TIME\_\_ & 编译时间，为四位整数常量 hhmm，例如 0014 或 2300 \\
\end{longtable}

例如，\cmd{\_\_PXL\_DATE\_\_} 与 \cmd{\_\_PXL\_TIME\_\_} 可用于编写基于当前日期改变行为的代码。因为这些变量是预处理器变量，也可用于 \cmd{\#if} 与 \cmd{\#define} 语句。
\begin{lstlisting}[style=pxl]
/**
 * define a new token, expire_20061231, that is be replaced
 * by the token deprecated before 2006.12.31 but expands
 * to obsolete thereafter.
 */
#if __PXL_DATE__ < 20061231
#define expire_20061231 deprecated
#endif

#ifndef expire_20061231
#define expire_20061231 obsolete
#endif
\end{lstlisting}

\subsubsection{发布版本宏}
\subsubsection*{Release Version Macros}

已发布宏允许你将当前发布版本与其他发布版本比较。

例如，你可用年、月、service pack 与 hot fix 的整数值判断当前发布版本是否至少与指定发布版本一样新。
\begin{lstlisting}[style=pxl]
#if VERSION_GE(2009, 6, 1, 0)
. . .      // code that uses new features introduced in 2009.06.SP01
#else
. . .      // code that uses older method (pre 2006.06.SP01)
#endif
\end{lstlisting}

VERSION 宏定义见表~\ref{tab:pxl-version-macros}。

\begin{noteBox}
\cmd{year}、\cmd{month}、\cmd{service\_pack} 与 \cmd{patch} 为整数值。
\end{noteBox}

\begin{longtable}{@{}>{\raggedright\arraybackslash\ttfamily}p{0.42\textwidth} p{0.50\textwidth}@{}}
\caption{VERSION 宏}\label{tab:pxl-version-macros}\\
\toprule
宏 & 定义 \\
\midrule
\endfirsthead
\caption[]{VERSION 宏（续）}\\
\toprule
宏 & 定义 \\
\midrule
\endhead
\bottomrule
\endfoot
VERSION\_GE(year, month, service\_pack, patch) & 测试当前发布版本是否大于或等于指定发布版本。 \\
VERSION\_LE(year, month, service\_pack, patch) & 测试当前发布版本是否小于或等于指定发布版本。 \\
VERSION\_GT(year, month, service\_pack, patch) & 测试当前发布版本是否大于指定发布版本。 \\
VERSION\_LT(year, month, service\_pack, patch) & 测试当前发布版本是否小于指定发布版本。 \\
VERSION\_EQ(year, month, service\_pack, patch) & 测试当前发布版本是否等于指定发布版本。 \\
\end{longtable}

\subsubsection{日期比较宏}
\subsubsection*{Date Comparison Macros}

日期比较宏允许你根据当前日期改变 runset 行为。

例如，可用新宏按当前日期启用或禁用某些 PXL 代码。
\begin{lstlisting}[style=pxl]
if (DATE_EQ(2015, 6, 4)) {
    note("Today == 2015-06-04");
} else {
    note("Today != 2015-06-04");
}
\end{lstlisting}

DATE 宏定义见表~\ref{tab:pxl-date-macros}。

\begin{noteBox}
\cmd{year}、\cmd{month} 与 \cmd{day} 为整数值。指定无效或越界值（例如 day 值大于 31，或 month 值大于 12）会导致未定义行为。
\end{noteBox}

\begin{longtable}{@{}>{\raggedright\arraybackslash\ttfamily}p{0.36\textwidth} p{0.56\textwidth}@{}}
\caption{DATE 宏}\label{tab:pxl-date-macros}\\
\toprule
宏 & 定义 \\
\midrule
\endfirsthead
\caption[]{DATE 宏（续）}\\
\toprule
宏 & 定义 \\
\midrule
\endhead
\bottomrule
\endfoot
DATE\_EQ(year, month, day) & 测试当前日期是否等于由 year、month、day 整数指定的日期。 \\
DATE\_GE(year, month, day) & 测试当前日期是否大于或等于由 year、month、day 整数指定的日期。 \\
DATE\_GT(year, month, day) & 测试当前日期是否大于由 year、month、day 整数指定的日期。 \\
DATE\_LE(year, month, day) & 测试当前日期是否小于或等于由 year、month、day 整数指定的日期。 \\
DATE\_LT(year, month, day) & 测试当前日期是否小于由 year、month、day 整数指定的日期。 \\
\end{longtable}

\subsection{限定符}
\subsubsection*{Qualifiers}

限定符通过改变其所标记项的行为或允许的操作来影响该项。表~\ref{tab:pxl-qualifiers} 定义各限定符。

\begin{longtable}{@{}>{\raggedright\arraybackslash\ttfamily}p{0.22\textwidth} p{0.70\textwidth}@{}}
\caption{限定符}\label{tab:pxl-qualifiers}\\
\toprule
限定符 & 定义 \\
\midrule
\endfirsthead
\caption[]{限定符（续）}\\
\toprule
限定符 & 定义 \\
\midrule
\endhead
\bottomrule
\endfoot
const & 每次对该对象赋新值都会导致错误。初始化不导致此错误。仅在尝试直接赋值或通过带 \cmd{out} 或 \cmd{in\_out} 限定的函数参数传递时才会报错。 \\
deprecated & 每次定义或引用该对象都会导致警告。在函数原型中为函数参数赋初值，或在该参数所属函数体内使用该参数，不会生成此警告。 \\
in\_out & 函数参数在每次函数调用中要求可赋值项（lvalue）。不可赋值项或类型错误的项会导致错误。该项的值也会因函数调用而改变。关于 lvalue 的更多信息，见「赋值运算符」（原书第~199~页）。 \\
inline & 将函数标记为在使用 \cmd{-ndg} 命令行选项的 IC Validator 运行中需要展开。见 IC Validator 命令行选项（原书第~25~页）。 \\
literal & 函数参数要求函数实参为字面量表达式。也就是说，编译器必须能在解析时解析函数实参的值。 \\
obsolete & 每次定义或引用该项都会导致错误。在函数原型中为函数参数赋初值不会生成此错误。 \\
out & 函数参数在每次函数调用中要求可赋值项（lvalue）。不可赋值项或类型错误的项会导致错误。处理函数体时假定该参数无值。该项的值也会因函数调用而改变。关于 lvalue 的更多信息，见「赋值运算符」（原书第~199~页）。 \\
\end{longtable}

限定符只能用于限定相关项。限定符见表~\ref{tab:pxl-qualifier-items}。

\begin{longtable}{@{}>{\raggedright\arraybackslash\ttfamily}p{0.20\textwidth}
  >{\centering\arraybackslash}p{0.16\textwidth}
  >{\centering\arraybackslash}p{0.16\textwidth}
  >{\centering\arraybackslash}p{0.18\textwidth}
  >{\centering\arraybackslash}p{0.18\textwidth}@{}}
\caption{限定符适用项}\label{tab:pxl-qualifier-items}\\
\toprule
限定符 & 变量 & 函数 & 参数 & 结构成员 \\
\midrule
\endfirsthead
\caption[]{限定符适用项（续）}\\
\toprule
限定符 & 变量 & 函数 & 参数 & 结构成员 \\
\midrule
\endhead
\bottomrule
\endfoot
binary & & 是 & & \\
const & 是 & & 是 & \\
deprecated & 是 & 是 & 是 & 是 \\
in\_out & & & 是 & \\
inline & & 是 & & \\
literal & & & 是 & \\
obsolete & 是 & 是 & 是 & 是 \\
out & & & 是 & \\
\end{longtable}

% ============================================================
\section{模块化代码示例}
\subsection*{Example of Modularized Code}

\begin{itemize}
  \item \cmd{function.rh}
\begin{lstlisting}[style=pxl]
// Use pragma once to prevent including
// this file multiple times
#pragma once

#include <icv.rh>

// Function prototypes
fx1 : function (
   lyr : polygon_layer,
   w   : double
) returning out: polygon_layer;

fx2 : function (
   lyr : polygon_layer,
   d   : double
) returning out: polygon_layer;

#include <functions.rs>
\end{lstlisting}

  \item \cmd{function.rs}
\begin{lstlisting}[style=pxl]
// Function definitions
fx1 : function (
   lyr : polygon_layer,
   w   : double
   ) returning out: polygon_layer {
      // code body
      out = f( lyr, w);
   };

fx2 : function (
   lyr : polygon_layer,
   d   : double
   ) returning out: polygon_layer {
      // code body
      out = f( lyr, d);
   };
\end{lstlisting}

  \item \cmd{main.rs}
\begin{lstlisting}[style=pxl]
#include <icv.rh>
#include <function.rh>

// assign(), library(), etc.

y = fx1( a, 0.1);
z = fx2( b, 0.3);
\end{lstlisting}
\end{itemize}

